\documentclass[11pt,twoside]{article}

\usepackage[margin=1in]{geometry}
\usepackage{amsmath,amssymb,amsthm,mathtools}
\usepackage{array,booktabs}
\usepackage[inline,shortlabels]{enumitem}
\usepackage{xcolor}
\usepackage{microtype}
\usepackage{tikz}
\usetikzlibrary{arrows.meta,positioning,calc,fit}
\usepackage{../exercises/course-exercises}

% Revision date: maintained by scripts/build_lecture_notes.py. Keep this command.
\newcommand{\lecturerevised}{2026-10-25 23:41 UTC-06:00}
\usepackage{eso-pic}
\AddToShipoutPictureBG{%
  \AtPageLowerLeft{%
    \put(\LenToUnit{0.5\paperwidth},\LenToUnit{0.5in}){%
      \makebox(0,0){\normalfont\footnotesize\color{black!60}Last revised: \lecturerevised}%
    }%
  }%
}

\usepackage[square,authoryear]{natbib}
\usepackage[colorlinks=true,linkcolor=blue!60!black,
  urlcolor=blue!60!black,citecolor=blue!60!black]{hyperref}
\usepackage[nameinlink,capitalize]{cleveref}
\crefname{exercise}{Exercise}{Exercises}
\Crefname{exercise}{Exercise}{Exercises}
\newcommand{\appendixref}[1]{\hyperref[#1]{Appendix~\ref*{#1}}}

\usepackage{enotez}
\DeclareInstance{enotez-list}{lecture}{list}{format=\normalsize,number={#1.}}
\setenotez{list-name=Endnotes,list-style=lecture,backref=true}
\crefname{endnote}{Endnote}{Endnotes}
\Crefname{endnote}{Endnote}{Endnotes}

\definecolor{courseblue}{HTML}{17365D}
\definecolor{lightblue}{HTML}{DCE6F1}
\newcommand{\E}{\mathbb E}
\newcommand{\I}{\mathbb I}
\renewcommand{\P}{\mathbb P}
\newcommand{\Law}{\mathcal L}
\newcommand{\R}{\mathbb R}
\newcommand{\relu}{\operatorname{ReLU}}
\newcommand{\TV}{\operatorname{TV}}
\makeatletter
\@for\@letter:=A,B,C,D,E,F,G,H,I,J,K,L,M,N,O,P,Q,R,S,T,U,V,W,X,Y,Z\do{%
  \expandafter\edef\csname c\@letter\endcsname{\noexpand\mathcal{\@letter}}}
\makeatother
\newcommand{\lecturenumber}{15}

\newtheorem{theorem}{Theorem}[section]
\newtheorem{proposition}[theorem]{Proposition}
\newtheorem{corollary}[theorem]{Corollary}
\newtheorem{lemma}[theorem]{Lemma}
\theoremstyle{definition}
\newtheorem{definition}[theorem]{Definition}
\newtheorem{example}[theorem]{Example}
\theoremstyle{remark}
\newtheorem{remark}[theorem]{Remark}

\renewcommand{\thepage}{15--\arabic{page}}
\renewcommand{\thesection}{15.\arabic{section}}
\renewcommand{\theequation}{15.\arabic{equation}}
\pagestyle{myheadings}
\markboth{Lecture 15: When neural networks approximate efficiently}
{Lecture 15: When neural networks approximate efficiently}

\begin{document}

\begin{center}
\fcolorbox{courseblue}{white}{%
\begin{minipage}{0.92\textwidth}
\vspace{1mm}
\textbf{CMPUT 654: Mathematical Foundations of Modern AI Systems}
\hfill Fall 2026

\vspace{3mm}
\begin{center}
{\Large\bfseries Lecture 15: When neural networks approximate efficiently}

\vspace{1mm}
October 22, 2026
\end{center}

\vspace{2mm}
\textit{Lecturer: Csaba Szepesv\'ari}
\hfill
\textit{Note prepared by: Csaba Szepesv\'ari}
\vspace{1mm}
\end{minipage}}
\end{center}

\noindent\textbf{Reading guide.}
\href{lecture14.pdf}{Lecture 14} studied representation questions for maps with discrete inputs and
outputs. Lecture 15 changes both the domain and the target: the input belongs
to a subset of $\R^d$, the target is real-valued, and the statistical problem
is regression. Continuous regression makes three questions distinct. An
approximation result asks whether a good member of a function family exists.
A constructive approximation result asks whether an algorithm can produce one
from specified information about the target. A learning result asks what can
be estimated from noisy data.

The material presented in class runs through \cref{sec:mousavi}. The
development begins with the progression from approximation to constructive
approximation to learning.
Shallow ridge and radial-basis-function networks make the questions concrete.
The minimax limit for ordinary smooth regression supplies an
information-theoretic benchmark. The central Hilbert-space lemma then
shows that a function with a bounded integral representation has an accurate
finite approximation. The lemma's proof uses random sampling, although the
lemma's conclusion is deterministic. The two density corollaries answer the
minimum approximation question: radial-basis-function and ridge networks can
make the approximation error arbitrarily small in the stated \(L^2\) spaces.
Quantitative approximation results for H\"older balls describe how the error
decreases with the number of terms, and suitable neural-network constructions
give corresponding rates. The integral representation also identifies
functions that are easy to approximate by \(m\)-term networks. Barron's
theorem shows that the same class has the following property: for every fixed
\(m\)-dimensional linear space, some function in the class is poorly
approximated by that space, even though target-dependent ridge atoms attain an
\(m^{-1/2}\) approximation rate.

The Barron comparison makes target-dependent dictionary selection the central
issue. Fixing the dictionary distribution in advance gives a kernel method.
Allowing the dictionary elements to depend on the target creates a
feature-selection problem: can a constructive or statistical method find the
useful elements? Projection pursuit gives a greedy constructive approach
conditional on a difficult oracle. The general learning question remains
open, so the final two main sections restrict the target to a multi-index
model. Under a rank condition, target gradients and a singular-value
decomposition identify a predictive subspace. A special two-stage
stochastic-gradient procedure then gives a learning result under restrictive
assumptions.

The optional-reading section contains the entropy result that was not
presented in class, additional classical and neural comparators, and results
about depth.
The appendices contain the wavelet, Fourier, active-query, RKHS, and
staged-SGD details used to support the main development.
\href{lecture09.pdf}{Lecture 9},
\href{lecture09.pdf\#nameddest=appendix.A}{Appendix~9.A}, reviews $L^2(P_X)$,
inner products, Hilbert spaces, orthogonality, and projections.

The lecture's central question is:
\begin{center}
\emph{Neural networks can choose features for the target. Can a constructive
or statistical algorithm find the useful features from the information
available to the algorithm?}
\end{center}
The lecture gives a general representation theorem and a restricted learning
theorem. A general efficient learning theorem remains open.

\section{From approximation to regression learning}
\label{sec:questions}

Regression starts from random input--response pairs $(X,Y)$ with
$X\in\cX\subseteq\R^d$ and $Y\in\R$. Under squared loss, the target is the
conditional mean
\[
 f(x)=\E[Y\mid X=x].
\]
The input domain and the target are continuous. Regression leads to a
progression from approximation, to constructive approximation, to learning.

\paragraph{Approximation.}
For a target $f$ and a family $\cF_m$, how small can
\begin{equation}
 \inf_{g\in\cF_m}\|f-g\|
 \label{eq:approximation-error}
\end{equation}
be? The infimum is an existential quantity and may use complete knowledge of
$f$.

To make the approximation question concrete, consider the function families
used throughout the lecture. For an activation $\sigma:\R\to\R$, a width-$m$
one-hidden-layer \emph{ridge network} belongs to
\begin{equation}
 \cN_m^\sigma
 =\left\{x\mapsto c_0+\sum_{j=1}^m
 c_j\sigma(a_j^\top x+b_j):
 a_j\in\R^d,\ b_j,c_j,c_0\in\R\right\}.
 \label{eq:ridge-network-class}
\end{equation}
For $\psi:\R^d\to\R$, a width-$m$ scale--location network belongs to
\begin{equation}
 \cR_m^\psi
 =\left\{x\mapsto c_0+\sum_{j=1}^m
 c_j\psi\bigl(s_j(x-u_j)\bigr):
 s_j>0,\ u_j\in\R^d,\ c_j,c_0\in\R\right\}.
 \label{eq:scale-location-network-class}
\end{equation}
If $\psi(z)=\rho(\|z\|_2)$, then
$\psi(s_j(x-u_j))=\rho(s_j\|x-u_j\|_2)$, and
\cref{eq:scale-location-network-class} is a \emph{radial-basis-function
network}. The center $u_j$ selects a location, and the scale $s_j$ selects the
spatial resolution of the $j$th atom. These formulas specify representation
families. An approximation theorem may use complete knowledge of the target
to choose the parameters of each network.

\paragraph{Constructive approximation.}
Given an information oracle for $f$, can a finite algorithm call the oracle
and produce an accurate $g$? A point-evaluation oracle returns
\[
 L_x(f)=f(x)=\int f\,d\delta_x,
\]
where $\delta_x$ is the Dirac point mass. A general linear-information oracle
returns $L(f)$ for specified linear functionals $L$. Point evaluations are a
special case. A numerically stable construction changes little when the oracle
values are perturbed. Numerical stability permits a subsequent analysis of
observation noise.

\paragraph{Learning.}
Regression supplies noisy point evaluations. From random noisy observations,
can an estimator $\widehat f_n$ achieve small prediction risk? For the
$L^2(P_X)$ loss, the risk is
\begin{equation}
 \E\bigl[\|\widehat f_n-f\|_{L^2(P_X)}^2\bigr].
 \label{eq:learning-risk}
\end{equation}
The expectation averages over the training sample and any randomness used by
the estimator.

The norm used in the main results is
\begin{equation}
 \langle g,h\rangle_{L^2(P_X)}=\E[g(X)h(X)],
 \qquad
 \|h\|_{L^2(P_X)}^2=\E[h(X)^2],
 \label{eq:l2-inner-product}
\end{equation}
where $\Law(X)=P_X$. Functions that agree $P_X$-almost everywhere represent
the same element of $L^2(P_X)$. The space is a Hilbert space, so independent
centered errors add through the inner product exactly as Euclidean vectors do.
\href{lecture09.pdf}{Lecture 9},
\href{lecture09.pdf\#nameddest=appendix.A}{Appendix~9.A}, gives the required
background. Approximation theory also commonly uses the uniform norm
$\|h\|_\infty=\sup_x|h(x)|$; the present
lecture uses $L^2(P_X)$ because squared-loss regression averages prediction
error under the input law.\endnote{Uniform error is the appropriate objective
when every input must receive a guarantee. Such bounds also arise in dynamic
programming when a worst-case value-function error must remain controlled
through repeated Bellman updates. The sampling lemma below extends to any
Hilbert space, but $L^\infty$ is not a Hilbert space, so the same variance
calculation does not directly give an $L^\infty$ bound.}

No choice of representation or algorithm can avoid the information-theoretic
limits of regression. To state the standard benchmark, set
$\cX=[0,1]^d$. Suppose that $P_X$ has a
density bounded above and bounded away from zero, and let
$(X_i,Y_i)_{i=1}^n$ be independent copies of $(X,Y)$ satisfying
\begin{equation}
 Y_i=f(X_i)+\varepsilon_i,
 \qquad
 \Law(\varepsilon_i)=N(0,\sigma^2),
 \label{eq:regression-model}
\end{equation}
where the noise variables are independent of the inputs. A function $h$ is
\emph{$\beta$-H\"older} for $0<\beta\leq1$ if
\begin{equation}
 |h(x)-h(x')|\leq L\|x-x'\|_2^\beta
 \quad\text{for every }x,x'.
 \label{eq:beta-holder}
\end{equation}
The case $\beta=1$ is Lipschitz continuity. When $0<\beta<1$, small distances
satisfy $r^\beta\geq r$ for $0\leq r\leq1$, so the same constant permits
larger local changes than a Lipschitz condition.

Write $\alpha=k+\beta$ with $k\in\mathbb N_0$ and $0<\beta\leq1$. The
H\"older ball $\cH_d^\alpha(L)$ contains functions whose partial derivatives
through order $k$ are bounded by $L$ and whose derivatives of order $k$ are
$\beta$-H\"older with constant $L$. The minimax mean integrated squared error
over \(\cH_d^\alpha(L)\)
satisfies
\begin{equation}
 \inf_{\widehat f_n}\sup_{f\in\cH_d^\alpha(L)}
 \E\bigl[\|\widehat f_n-f\|_{L^2(P_X)}^2\bigr]
 \asymp n^{-2\alpha/(2\alpha+d)}.
 \label{eq:generic-smooth-rate}
\end{equation}
The infimum ranges over all measurable estimators based on the $n$ observed
pairs; the notation does not impose a computational model. Thus
\cref{eq:generic-smooth-rate} is an information-theoretic statement, not a
claim about an efficiently implementable algorithm. The exponent deteriorates
with the ambient dimension. A faster rate requires a smaller target class, a
different observation model, or a different loss. The bound sets the
regression benchmark. A complementary entropy result was not developed in
class and appears in the optional reading.

\section{Integral representations and density}
\label{sec:sampling}

The finite network formulas in
\cref{eq:ridge-network-class,eq:scale-location-network-class} both start from
a \emph{dictionary}: a parameterized family
\(\cD=(\phi_\theta)_{\theta\in\Theta}\) of functions. For a ridge network,
\(\theta=(a,b)\) and
\(\phi_\theta(x)=\sigma(a^\top x+b)\). For a scale--location network,
\(\theta=(s,u)\) and \(\phi_\theta(x)=\psi(s(x-u))\). A finite network
selects finitely many dictionary elements. Replacing the finite sum by an
integral permits a continuous collection and gives a useful sufficient
condition for accurate finite approximation.

\begin{definition}[Finite signed measure and integral representation]
\label{def:signed-integral-representation}
Let \(\Theta\) be a measurable parameter space. A
\emph{finite signed measure} will mean a measure of the form
\[
 d\mu(\theta)=s(\theta)\,d\nu(\theta),
\]
where \(\nu\) is a finite positive measure and
\(s:\Theta\to\{-1,1\}\) is measurable. The positive measure
\(|\mu|:=\nu\) is the \emph{total variation measure} of \(\mu\), and
\[
 \|\mu\|_{\TV}:=|\mu|(\Theta)
\]
is its \emph{total variation norm}. For a dictionary
\((\phi_\theta)_{\theta\in\Theta}\) in a Hilbert space \(\cH\), an
\emph{integral representation} of \(f\in\cH\) is
\begin{equation}
 f=\int_\Theta \phi_\theta\,d\mu(\theta).
 \label{eq:integral-representation}
\end{equation}
\end{definition}
The definition through a sign function is equivalent to the usual definition
of a finite real-valued signed measure.%
\endnote{The usual definition starts from a countably additive real-valued
set function. The Jordan decomposition theorem writes every such measure as
\(\mu=\mu^+-\mu^-\) for mutually singular positive measures. Consequently,
\(|\mu|=\mu^++\mu^-\) and
\(d\mu=s\,d|\mu|\) for a measurable sign function
\(s\in\{-1,1\}\), up to a \(|\mu|\)-null set. Thus the sign-times-positive-
measure definition loses no finite signed measures.}

A finite sum is a special integral. If
\(f=\sum_{j=1}^m c_j\phi_{\theta_j}\), take
\(\mu=\sum_{j=1}^m c_j\delta_{\theta_j}\). Then
\(|\mu|=\sum_{j=1}^m|c_j|\delta_{\theta_j}\) after repeated atoms have been
combined. Integral representations can also describe functions that require
a limit of finite sums.

The same formula is used in mean-field descriptions of two-layer networks. If
\(\rho\) is a probability measure on neuron parameters, then
\begin{equation}
 f_\rho(x)=\int a\,\sigma(w^\top x+b)\,d\rho(a,w,b).
 \label{eq:mean-field-representation}
\end{equation}
Absorbing the output weight \(a\) into a signed measure on \((w,b)\) gives
the representation in \cref{eq:integral-representation}. Mean-field training
adds an evolution equation for \(\rho\); the static integral formula alone is
a representation statement.

The integral form is useful because every function with a representing
measure of bounded total variation has an accurate finite-sum approximation.
The following lemma is the central result of the lecture. The lemma's statement
requires only Hilbert-space geometry.

\begin{lemma}[Finite approximation of an integral representation]
\label{lem:mc-sparsification}
Let \(\cH\) be a real Hilbert space, let
\(\theta\mapsto\phi_\theta\in\cH\) be measurable, and suppose that
\(\|\phi_\theta\|_{\cH}\leq1\) for every \(\theta\). If
\(f=\int\phi_\theta\,d\mu(\theta)\) and
\(B=\|\mu\|_{\TV}<\infty\), then, for every \(m\geq1\), there are
\(\theta_1,\ldots,\theta_m\in\Theta\) and
\(s_1,\ldots,s_m\in\{-1,1\}\) such that
\begin{equation}
 \left\|f-\frac{B}{m}\sum_{j=1}^m
 s_j\phi_{\theta_j}\right\|_{\cH}
 \leq\frac{B}{\sqrt m}.
 \label{eq:mc-sparsification-rate}
\end{equation}
\end{lemma}

\begin{proof}
The conclusion is immediate when \(B=0\). Suppose that \(B>0\). Let
\(\Theta_1,\ldots,\Theta_m\) be independent random elements with common law
\(\pi=|\mu|/B\), let \(s=d\mu/d|\mu|\), and define
\[
 S_j=s(\Theta_j),
 \qquad
 Z_j=B S_j\phi_{\Theta_j}.
\]
The random elements \(Z_1,\ldots,Z_m\) are independent,
\(\E[Z_j]=f\), and \(\|Z_j\|_{\cH}\leq B\). The inner product and
independence give
\begin{align}
 \E\left[
 \left\|\frac1m\sum_{j=1}^mZ_j-f\right\|_{\cH}^2\right]
 &=\frac1m\E\bigl[\|Z_1-f\|_{\cH}^2\bigr]
 =\frac1m\left(\E[\|Z_1\|_{\cH}^2]-\|f\|_{\cH}^2\right)
 \leq\frac{B^2}{m}.
 \label{eq:mc-proof}
\end{align}
At least one realization has squared error no larger than the expectation.
Fixing that realization proves \cref{eq:mc-sparsification-rate}.
\end{proof}

The proof is the \emph{probabilistic method}: a random experiment proves that
a deterministic choice exists. The proof also exposes the limitation of the
construction. The sampling law \(|\mu|/B\) depends on a representing measure
for \(f\). Knowing that measure can require essentially the same information
as knowing the target. The lemma therefore proves economical finite
representation. The lemma alone does not supply an implementable oracle
algorithm or a regression estimator.

The lemma gives a route from integral representations to universal
approximation. If a dense subclass of targets has representing measures of
finite total variation, finite networks are dense as well. Applying the
integral-representation argument to scale--location atoms and ridge atoms gives the next two
corollaries.

For a measurable domain \(D\subseteq\R^d\), \(L^2(D)\) denotes the space of
measurable functions \(h:D\to\R\)
such that \(\int_D|h(x)|^2\,dx<\infty\), where functions that agree
Lebesgue-almost everywhere are identified. The \(L^2(D)\) norm is
\(\|h\|_{L^2(D)}=(\int_D|h(x)|^2\,dx)^{1/2}\). For a subset
\(\cF\subseteq L^2(D)\), the notation
\(\overline{\cF}^{\,L^2(D)}\) denotes the closure of \(\cF\) in the
\(L^2(D)\) norm. Thus
\(\overline{\cF}^{\,L^2(D)}=L^2(D)\) means that for every
\(h\in L^2(D)\) and every \(\epsilon>0\), there is an \(f\in\cF\) such
that \(\|h-f\|_{L^2(D)}<\epsilon\).

\begin{corollary}[Density of Gaussian radial-basis-function networks]
\label{cor:rbf-density}
Let \(\psi(x)=\exp(-\|x\|_2^2/2)\). Then
\begin{equation}
 \overline{\bigcup_{m\geq1}
 \left\{\sum_{j=1}^m c_j
 \psi\bigl(s_j(\cdot-u_j)\bigr):
 c_j\in\R,\ s_j>0,\ u_j\in\R^d\right\}}^{\,L^2(\R^d)}
 =L^2(\R^d).
 \label{eq:rbf-density}
\end{equation}
\end{corollary}

A wavelet proof chooses a Laplacian-of-Gaussian analysis function and the
Gaussian as the synthesis atom. Translation and dilation then give an
integral representation, and \cref{lem:mc-sparsification} replaces the
integral by finitely many scale--location atoms.
\appendixref{app:wavelet-facts} and \appendixref{app:fourier-basics} give the
coefficient estimate, reconstruction formula, and Fourier verification.

\begin{corollary}[Density of shallow ridge networks on a cube]
\label{cor:ridge-density}
Let \(Q=[0,1]^d\), and let \(\sigma:\R\to[0,1]\) be measurable with
\begin{equation}
 \lim_{t\to-\infty}\sigma(t)=0,
 \qquad
 \lim_{t\to+\infty}\sigma(t)=1.
 \label{eq:sigmoidal-activation}
\end{equation}
For
\begin{equation}
 \cN_\sigma(Q)
 =\left\{\left.
 c_0+\sum_{j=1}^m c_j\sigma(a_j^\top x+b_j)
 \right|_Q:
 m\geq1,\ c_0,c_j,b_j\in\R,\ a_j\in\R^d\right\},
 \label{eq:finite-ridge-class}
\end{equation}
Then
\begin{equation}
 \overline{\cN_\sigma(Q)}^{\,L^2(Q)}=L^2(Q).
 \label{eq:ridge-density}
\end{equation}
\end{corollary}

Fourier inversion writes a smooth target as an integral of oscillations
\(e^{i\omega^\top x}\). On the cube, each oscillation is a univariate
function of \(\omega^\top x\), and the fundamental theorem of calculus writes
the univariate profile as an integral of thresholds. Steep sigmoids
approximate the thresholds. \appendixref{app:fourier-basics} makes each step
precise.

The two density corollaries answer a basic minimum question: can a finite
network make the approximation error arbitrarily small? The answer is yes for
the stated RBF and ridge families. Density alone gives no relation between the
required network width and the desired error.

For a H\"older ball, approximation theory asks the quantitative version of
the same question. Under the input-distribution assumptions used in
\cref{eq:generic-smooth-rate}, standard local-polynomial or wavelet spaces
\(V_m\) with \(m\) degrees of freedom can be chosen so that
\begin{equation}
 \sup_{f\in\cH_d^\alpha(L)}\inf_{g\in V_m}
 \|f-g\|_{L^2(P_X)}\lesssim m^{-\alpha/d}.
 \label{eq:holder-approximation-rate}
\end{equation}
Suitable neural-network architectures can implement the same local
approximations. For example, ReLU networks with \(O(m\log m)\) parameters can
attain uniform, and hence \(L^2(P_X)\), error of order
\(m^{-\alpha/d}\) by implementing local Taylor approximations and a partition
of unity. The parameter count and logarithmic factors depend on the activation
and architecture. For a classical \(m\)-dimensional series estimator, the
squared approximation error is of order \(m^{-2\alpha/d}\) by
\cref{eq:holder-approximation-rate}, and the estimation error is of order
\(m/n\). Balancing the two terms gives the regression rate
\(n^{-2\alpha/(2\alpha+d)}\) in \cref{eq:generic-smooth-rate}. The bias--
estimation balance connects the quantitative approximation result to the
regression benchmark; approximation and regression still answer different
questions.

\section{Target-dependent approximation and fixed features}
\label{sec:adaptive-fixed-features}

The density results are useful, but there is a more interesting fact. Some
functions have a bounded continuous neural-network representation of the form
in \cref{eq:integral-representation}. Lemma~\ref{lem:mc-sparsification}
therefore gives accurate \(m\)-term approximations for those functions. Now fix
any \(m\) basis functions in advance, whether they come from a classical
approximation method or from any other proposed basis. Barron's theorem says
that some function from the same class is nevertheless poorly approximated by
their span. To state the theorem, first define the target class by fixing the
Fourier normalization
\begin{equation}
 \widehat f(\omega)=\int_{\R^d}e^{-i\omega^\top x}f(x)\,dx.
 \label{eq:fourier-transform}
\end{equation}

\begin{definition}[Anchored spectral Barron norm]
\label{def:spectral-barron-norm}
Let \(\cB(\R^d)\) be the class of real-valued functions
\(f:\R^d\to\R\) that admit a representation
\begin{equation}
 f(x)=c+\frac1{(2\pi)^d}\int_{\R^d}
 (e^{i\omega^\top x}-1)F(\omega)\,d\omega,
 \label{eq:barron-density-representation}
\end{equation}
where \(c\in\R\), \(F\in L^1(\R^d;\mathbb C)\), the integral is real for
every \(x\), and
\(\int_{\R^d}\|\omega\|_1|F(\omega)|\,d\omega<\infty\). Setting \(x=0\) in
\cref{eq:barron-density-representation} gives \(c=f(0)\). The weighted spectral
term assigns value zero to every constant function by taking \(F=0\). Using
only the weighted spectral term would therefore give a seminorm. Adding
\(|c|\) records the constant component. Define
\begin{equation}
 \|f\|_{\cB(\R^d)}
 =\inf_{c,F}\left\{
 |c|+\frac1{(2\pi)^d}\int_{\R^d}
 \|\omega\|_1|F(\omega)|\,d\omega\right\},
 \label{eq:barron-norm}
\end{equation}
where the infimum ranges over all representations in
\cref{eq:barron-density-representation}. The ambient Barron ball and its
restriction to \(Q=[0,1]^d\) are
\begin{equation}
 \cB_C(\R^d)=\{f\in\cB(\R^d):\|f\|_{\cB(\R^d)}\leq C\},
 \qquad
 \cB_C(Q)=\{f|_Q:f\in\cB_C(\R^d)\}\subseteq L^2(Q).
 \label{eq:barron-class}
\end{equation}
\end{definition}
Each restriction in \(\cB_C(Q)\) is identified with its equivalence class in
\(L^2(Q)\). The quantity in \cref{eq:barron-norm} is a norm on the represented
functions on \(\R^d\).%
\endnote{One may instead use the spectral moment as a norm modulo constants.
The network represents \(c\) exactly through its output bias, so the
approximation rate depends only on the spectral moment. Different Fourier
normalizations change constants but not the approximation exponents.}

\begin{theorem}[Barron's adaptive upper bound and fixed-space lower bound]
\label{thm:barron}
Let \(P_X\) be uniform on \(Q=[0,1]^d\), and let \(\sigma\) satisfy
\cref{eq:sigmoidal-activation}. There is a universal constant
\(\kappa>0\) such that the following statements hold.
\begin{enumerate}[label=(\roman*),leftmargin=*]
\item For every \(h\in\cB_C(Q)\), every \(m\geq1\), and every
\(\eta>0\), there is
\(f_m\in\cN_m^\sigma\) satisfying
\begin{equation}
 \|h-f_m\|_{L^2(P_X)}\leq\frac{2(C+\eta)}{\sqrt m}.
 \label{eq:barron-upper}
\end{equation}
\item For every fixed collection \(h_1,\ldots,h_m\in L^2(P_X)\),
\begin{equation}
 \sup_{h\in\cB_C(Q)}
 \inf_{a\in\R^m}
 \left\|h-\sum_{j=1}^ma_jh_j\right\|_{L^2(P_X)}
 \geq \kappa\frac{C}{d}m^{-1/d}.
 \label{eq:barron-lower}
\end{equation}
\end{enumerate}
\end{theorem}

\paragraph{Upper-bound idea.}
Choose \(f\in\cB_C(\R^d)\) satisfying \(f|_Q=h\), and choose a representation
of \(f\) whose cost is at most \(C+\eta\). The
Fourier-to-threshold calculation used for
\cref{cor:ridge-density} represents \(f-c\) by a signed measure
over ridge atoms whose total variation is at most twice the spectral moment.
Lemma~\ref{lem:mc-sparsification} then selects \(m\) atoms. The constant
\(c\) is represented exactly. The selected ridge directions and biases depend
on the target.

\paragraph{Lower-bound idea.}
Choose \(2m\) distinct integer frequencies of size
\(O(dm^{1/d})\). The associated cosine functions are orthonormal in
\(L^2(Q)\), and their Barron norms have the same order as their frequencies.
For any subspace \(V\) of dimension at most \(m\), Bessel's inequality gives
\[
 \sum_{j=1}^{2m}\|P_Ve_j\|_{L^2(Q)}^2\leq m.
\]
At least one normalized cosine \(e_j\) therefore has
\[
 \inf_{v\in V}\|e_j-v\|_{L^2(Q)}^2
 =1-\|P_Ve_j\|_{L^2(Q)}^2\geq\frac12.
\]
Scaling that cosine into \(\cB_C(Q)\) gives
\cref{eq:barron-lower}. \appendixref{app:fourier-basics} handles the harmless
smoothing needed to use an integrable spectral density.

The lower bound is a statement about the \emph{Kolmogorov \(m\)-width}
\begin{equation}
 d_m(\cF;L^2(P_X))
 =\inf_{\substack{V\subseteq L^2(P_X)\\ \dim(V)\leq m}}
 \sup_{f\in\cF}\inf_{v\in V}\|f-v\|_{L^2(P_X)}.
 \label{eq:kolmogorov-width}
\end{equation}
Kolmogorov's question is whether one \(m\)-dimensional subspace can
approximate every target in \(\cF\). Barron's upper bound allows a different
set of \(m\) ridge atoms for each target. To reach error \(\epsilon\), the
upper bound uses \(m=O(C^2\epsilon^{-2})\), while a fixed linear space can
require \(m=\Omega((C/(d\epsilon))^d)\).%
\endnote{The comparison is stated for the Barron ball of radius \(C\). The
theorem does not assert that \(C\) stays bounded along a high-dimensional
sequence of targets.}

The special feature in Barron's upper bound is that the ridge directions and
biases are chosen after the target is known. The alternative is to choose the
dictionary elements in advance. In the integral representation in
\cref{eq:integral-representation}, choosing the dictionary in advance means
fixing a probability measure \(\tau\) on \(\Theta\) and permitting only
measures of the form \(d\mu=p\,d\tau\). Define the resulting function space and
its coefficient norm by
\begin{align}
 \cF_{2,\tau}
 &=\left\{f:f(x)=\int_\Theta p(\theta)\phi_\theta(x)\,d\tau(\theta)
              \text{ for some }p\in L^2(\tau)\right\},
 \notag\\
 \|f\|_{\cF_{2,\tau}}
 &=\inf\left\{\|p\|_{L^2(\tau)}:
       f(x)=\int_\Theta p(\theta)\phi_\theta(x)\,d\tau(\theta)\right\}.
 \label{eq:F2}
\end{align}
The feature locations are distributed according to \(\tau\); only the
coefficient function \(p\) changes with the target. The corresponding kernel
is
\begin{equation}
 k_\tau(x,x')
 =\int_\Theta\phi_\theta(x)\phi_\theta(x')\,d\tau(\theta).
 \label{eq:feature-kernel}
\end{equation}
The normed function space in \cref{eq:F2} is the reproducing-kernel Hilbert
space associated with \(k_\tau\). The infimum is needed because different
coefficient functions can represent the same function. Conversely, every
separable reproducing-kernel Hilbert space admits a feature representation of
the form in \cref{eq:F2}; \appendixref{app:rkhs-features} gives a direct construction.

The integral representation in \cref{eq:integral-representation} allows any
finite signed measure. Define the corresponding variation space and norm by
\begin{align}
 \cF_1
 &=\left\{f:f(x)=\int_\Theta\phi_\theta(x)\,d\mu(\theta)
             \text{ for some finite signed measure }\mu\right\},
 \notag\\
 \|f\|_{\cF_1}
 &=\inf\left\{\|\mu\|_{\TV}:
       f(x)=\int_\Theta\phi_\theta(x)\,d\mu(\theta)\right\}.
 \label{eq:F1}
\end{align}
The two unit balls are
\begin{equation}
 B_{\cF_1}=\{f:\|f\|_{\cF_1}\leq1\},
 \qquad
 B_{\cF_{2,\tau}}=\{f:\|f\|_{\cF_{2,\tau}}\leq1\}.
 \label{eq:F1-F2-unit-balls}
\end{equation}
For either \(\cG\in\{\cF_1,\cF_{2,\tau}\}\), every nonzero function
\(f\in\cG\) is a positive scalar multiple of its normalized version in the
unit ball:
\(f=\|f\|_{\cG}(f/\|f\|_{\cG})\). Thus the unit ball represents every
direction in the function space; the norm records the required scale.

The fixed-kernel representation requires \(d\mu=p\,d\tau\) with
\(p\in L^2(\tau)\). Since \(\tau\) is a probability measure,
Cauchy--Schwarz gives
\begin{equation}
 \|f\|_{\cF_1}
 \leq\int_\Theta|p|\,d\tau
 \leq\|p\|_{L^2(\tau)}.
 \label{eq:F2-F1-norm-inclusion}
\end{equation}
Taking the infimum over all coefficient functions representing \(f\) yields
\(\|f\|_{\cF_1}\leq\|f\|_{\cF_{2,\tau}}\), and hence
\begin{equation}
 \cF_{2,\tau}\subseteq\cF_1,
 \qquad B_{\cF_{2,\tau}}\subseteq B_{\cF_1}.
 \label{eq:F2-F1-ball-inclusion}
\end{equation}
Either inclusion can be strict, depending on the dictionary and feature law.
The mathematical distinction is that a variation-space representation may
concentrate mass on features chosen for the target, while an RKHS
representation must have an \(L^2(\tau)\) density relative to the fixed law
\(\tau\).

The representation results show a possible advantage of neural networks: the
dictionary elements can be chosen for the target instead of being fixed in
advance. Can a constructive or statistical algorithm find the useful
dictionary elements? Projection pursuit is a classical constructive attempt.

\section{Constructive approximation by projection pursuit}
\label{sec:structure}

\emph{Projection pursuit regression} selects target-dependent ridge functions
greedily. Projection pursuit approximates a target by
\begin{equation}
 f_K(x)=c_0+\sum_{j=1}^K g_j(a_j^\top x),
 \qquad a_j\in\R^d,
 \label{eq:projection-pursuit-model}
\end{equation}
where both the directions \(a_j\) and the univariate profiles \(g_j\) are
selected for the target. A shallow neural network is a restricted
projection-pursuit model with
\(g_j(t)=c_j\sigma(t+b_j)\).

A stagewise constructive method begins with a residual \(r_0=f\). At step
\(j\), the method searches for a direction and profile whose ridge function
has large correlation with \(r_{j-1}\), adds that function to the current
approximation, and updates the residual. Backfitting may then revise earlier
terms. In an abstract Hilbert space, an exact correlation oracle and a bounded
variation norm give an \(O(K^{-1/2})\) approximation bound.
\appendixref{app:active-query-recovery} states and proves the greedy result.

The oracle is the main computational issue. For a finite dictionary, one can
scan every atom. For the continuous ridge dictionary, the oracle must optimize
over a direction and a univariate function, and the search is generally
nonconvex.

Projection pursuit therefore gives a constructive approximation method only
when a strong atom-selection oracle is available. Learning requires an
efficient algorithm that finds useful features from noisy data, without
access to the target or to such an oracle. An efficient learning algorithm for
the general Barron class remains open. The next positive result restricts the
target to a multi-index model, where the features form one low-dimensional
linear subspace.

\section{Multi-index models and average derivatives}
\label{sec:agop}

The multi-index restriction replaces the search for arbitrary ridge sums by
the recovery of one predictive subspace. Several common structured models
have the form
\begin{equation}
 f(x)=\sum_{j=1}^K g_j(W_j^\top x),
 \qquad W_j\in\R^{d\times s_j}.
 \label{eq:structured-family}
\end{equation}
An additive model uses coordinate directions. A single-index model has
\(K=s_1=1\). Projection pursuit uses \(s_j=1\) and unknown directions.
A one-hidden-layer network is projection pursuit with a specified parametric
form for the profiles. A \emph{multi-index model} focuses on one term,
\begin{equation}
 f(x)=g(W^\top x),
 \qquad W\in\R^{d\times s},\quad W^\top W=I_s,\quad s\ll d.
 \label{eq:multi-index-model}
\end{equation}
Only the column space of \(W\) matters: an orthogonal change of basis inside
that space can be absorbed into the unknown link \(g\).

Differentiation reveals the subspace. The chain rule gives
\begin{equation}
 \nabla f(x)=W\nabla g(W^\top x),
 \label{eq:multi-index-gradient}
\end{equation}
so every gradient belongs to \(\operatorname{span}(W)\). If gradients are
available at points \(x_1,\ldots,x_q\), form
\begin{equation}
 G=\begin{bmatrix}
 \nabla f(x_1)&\cdots&\nabla f(x_q)
 \end{bmatrix}\in\R^{d\times q}.
 \label{eq:main-sampled-gradient-matrix}
\end{equation}
When \(\operatorname{rank}(G)=s\), the columns of \(G\) span
\(\operatorname{span}(W)\). For a singular-value decomposition
\begin{equation}
 G=\sum_{\ell=1}^s\sigma_\ell u_\ell v_\ell^\top,
 \qquad \sigma_s>0,
 \label{eq:gradient-svd}
\end{equation}
the vectors \(u_1,\ldots,u_s\) therefore form an orthonormal basis of the
predictive subspace. If target gradients were available, the feature-selection
problem would therefore reduce to a singular-value decomposition. Recovering
the subspace is sufficient because the link can be redefined in the recovered
coordinates.

The subspace gives a dimension reduction. If \(g\) has H\"older smoothness
\(\alpha\) and the subspace were known, the generic regression benchmark
would become
\begin{equation}
 n^{-2\alpha/(2\alpha+s)}
 \label{eq:intrinsic-smooth-rate}
\end{equation}
instead of the ambient-dimensional rate in
\cref{eq:generic-smooth-rate}.

Regression observations do not contain the target gradients. Methods that
estimate averages of derivatives from regression data are called
\emph{average derivative estimation} methods. A population matrix that can
expose every direction is the \emph{average gradient outer product}
\begin{equation}
 G_f=\E[\nabla f(X)\nabla f(X)^\top]
 =W\,\E[\nabla g(W^\top X)\nabla g(W^\top X)^\top]W^\top.
 \label{eq:agop}
\end{equation}
If the middle matrix is positive definite, the range of \(G_f\) equals
\(\operatorname{span}(W)\). A zero or nearly singular middle matrix means
that the available derivatives fail to reveal one or more active directions.

Point evaluations can estimate gradients by finite differences in a noiseless
active-query model; \appendixref{app:active-query-recovery} gives a
quantitative construction followed by the SVD argument. Passive noisy
regression is harder because the derivative information must first be
estimated. Classical two-stage methods estimate the
index space by local smoothing and derivative averages and then fit the link
in the estimated coordinates. The error of the subspace estimate must not
dominate the second-stage link error.

Gradients and a singular-value decomposition would recover the predictive
subspace and reduce the regression problem from dimension \(d\) to dimension
\(s\). Regression data do not provide the gradients directly. The next result
asks whether a two-layer network trained from data can exploit the same
structure. The positive theorem uses an unusual staged procedure under
restrictive model and algorithmic assumptions.

\section{A theorem for neural feature learning by stochastic gradient descent}
\label{sec:mousavi}

The positive result concerns a particular two-stage procedure, not ordinary
simultaneous training of both layers. The student predictor is the two-layer
network
\begin{equation}
\widehat y(x;W,a,b)=\sum_{j=1}^m a_j\sigma(w_j^\top x+b_j),
\qquad W\in\R^{m\times d}.
\label{eq:two-layer-student}
\end{equation}
$W$ is the first-layer weight matrix, $b$ contains the hidden-unit biases, and
$a$ contains the second-layer or output weights. In the first stage, only
$W$ is updated. Every online-SGD update receives a fresh independent pair
$(X_t,Y_t)$ and has the form
\begin{equation}
 W^{(t+1)}=(1-\eta_t\lambda)W^{(t)}
 -\eta_t\nabla_W\ell\bigl(
   \widehat y(X_t;W^{(t)},a^{(0)},b^{(0)}),Y_t\bigr).
 \label{eq:mousavi-first-stage}
\end{equation}
$a^{(0)}$ and $b^{(0)}$ remain fixed throughout the first stage.%
\endnote{\appendixref{app:sgd-contraction} gives a more general first-stage
result: online SGD contracts the component of \(W\) orthogonal to a teacher
subspace. The appendix also records the population-gradient identity used in
the analysis.}

The prediction algorithm then uses a separate second stage. After computing
$W^{(T)}$, the algorithm draws independent new biases satisfying
$\Law(b_j)=\operatorname{Unif}[-\Delta,\Delta]$, freezes both $W^{(T)}$ and $b$,
and updates only the output weights for $t=0,\ldots,T'-1$:
\begin{equation}
 a^{(t+1)}=(1-\eta'_t\lambda')a^{(t)}
 -\eta'_t\nabla_a\ell\bigl(
   \widehat y(X_{I_t};W^{(T)},a^{(t)},b),Y_{I_t}\bigr),
 \label{eq:mousavi-second-stage}
\end{equation}
where the indices $I_t$ are sampled independently and uniformly from
$\{0,\ldots,T-1\}$ at the second-stage updates. The second stage
therefore reuses the same $T$ observations rather than receiving fresh data;
the sample count is $T$, while the two stages together perform $T+T'$ SGD
updates. Conditional on $W^{(T)}$ and $b$, the second stage is regularized
linear prediction in $a$. The guarantee concerns the staged algorithm and
does not cover ordinary simultaneous training.
The prediction theorem uses a symmetric first-stage
initialization: all rows of $W^{(0)}$ equal the same vector drawn from
$N(0,I_d/d)$, all output weights equal the same positive value of order
$1/m$, and all initial biases equal the same positive constant. The
second-stage resampling of $b$ replaces the initial biases. For the prediction theorem,
$\sigma(z)=z_+$, $\tau\geq1$, and
\begin{equation}
 \ell_H(z)=
 \begin{cases}
  z^2/2,& |z|\leq1,\\
  |z|-1/2,& |z|>1.
 \end{cases}
 \qquad
 R_\tau(h)=\E\bigl[\ell_H(h(X)-Y)\wedge\tau\bigr].
 \label{eq:huber-truncated-risk}
\end{equation}
\begin{theorem}[Single-index prediction by staged SGD, informal form]
Suppose \(\Law(X)=N(0,I_d)\), and suppose the response follows the single-index
model
\begin{equation}
 Y=f(u^\top X)+\varepsilon,
 \qquad \|u\|_2=1,
 \label{eq:single-index-model}
\end{equation}
where the noise \(\varepsilon\) is independent of \(X\), and where $f$ is
monotone and twice differentiable, $|f(0)|<1$, the second derivative $f''$ has
at most polynomial growth, and \(\varepsilon\) is $1$-sub-Gaussian. The
algorithm uses the two stages in \cref{eq:mousavi-first-stage,eq:mousavi-second-stage}.
For a sufficiently large second-stage iteration count $T'$ and under the
specified initialization and tuning assumptions, its returned predictor
$\widehat y_{T,T',m}$ satisfies, with high probability,
\begin{equation}
 R_\tau(\widehat y_{T,T',m})-\E[\ell_H(\varepsilon)]
 \leq \widetilde C\left(
       \sqrt{\frac{d+\log(1/\delta)}{T}}
       +\sqrt{\frac{\log(1/\delta)}{m}}\right),
 \label{eq:mousavi-risk}
\end{equation}
where $R_\tau$ is truncated population Huber risk, $\ell_H$ is the Huber
loss, and $\widetilde C$ suppresses logarithmic factors and constants
determined by the fixed link, truncation level, and tuning parameters.
\end{theorem}
For a fixed monotone polynomial link and fixed target accuracy, with $m$ and
$T'$ chosen sufficiently large, the required number $T$ of distinct samples
is $\widetilde O(d)$. In contrast, learning a degree-$p$ target with a
rotationally invariant kernel can require $d^{\Omega(p)}$ samples.

The theorem gives a positive learning result for the two-stage procedure. The
nearly linear dependence on \(d\) requires Gaussian covariates, explicit
weight decay, online first-stage samples, bias resampling, and separate
training of the two layers. The theorem does not show that ordinary
simultaneous training adapts across the structural families in
\cref{eq:structured-family}.

\section{Optional reading}
\label{sec:optional}

The material through \cref{eq:mousavi-risk} was presented in class. The
optional sections compare other feature-adaptation mechanisms, state the
entropy result that was not presented in class, and summarize what additional
depth can change.

\subsection{Additional feature-adaptation mechanisms}

Classical structure-adaptive smoothers alternate local derivative estimation
with an update of the smoothing geometry. Under model-specific smoothness,
design, and identifiability conditions, such procedures can estimate a
multi-index subspace before fitting the low-dimensional link. These classical
procedures show that
data-dependent directions are not unique to neural networks.

Variation-ball empirical risk minimization treats the representing measure in
\cref{eq:F1} as the fitted object. Statistical bounds can adapt to
projection-pursuit targets with a small hidden dimension, but standard
conditional-gradient implementations require a generally nonconvex
neuron-selection oracle.

The \emph{recursive feature machine} repeatedly fits a kernel predictor,
computes the empirical average gradient outer product of that predictor, and
uses the matrix to update the kernel metric. The method makes feature-metric
adaptation explicit without training a neural network. The \emph{neural
feature ansatz} is the empirical relation
\begin{equation}
 W_i^\top W_i
 \ \propto\
 \frac1n\sum_{r=1}^n
 \nabla f_i(h_i(X_r))\nabla f_i(h_i(X_r))^\top
 \label{eq:nfa}
\end{equation}
between the Gram matrix of a trained layer and the empirical AGOP of the
remaining network. Experiments support the relation across several settings;
available proofs use more restricted training schemes. The ansatz is an
empirically supported hypothesis, not a general theorem about ordinary deep
training.

\begin{table}[ht]
\centering
\small
\begin{tabular}{@{}>{\raggedright\arraybackslash}p{0.20\textwidth}
  >{\raggedright\arraybackslash}p{0.26\textwidth}
  >{\raggedright\arraybackslash}p{0.25\textwidth}
  >{\raggedright\arraybackslash}p{0.20\textwidth}@{}}
\toprule
Method & What changes with the data? & Positive statement & Main limitation under the stated assumptions \\
\midrule
Fixed linear space or fixed RKHS
& Coefficients only
& Standard estimation and approximation guarantees
& No learned directions or feature metric \\
Projection pursuit and structure-adaptive smoothing
& Ridge directions, term count, or smoothing geometry
& Greedy oracle rates; active-query and passive-regression guarantees
& Specialized observations and models; continuous atom search is nonconvex \\
Variation-ball ERM
& Representing measure over neurons
& Structured statistical rate with hidden dimension $s$
& Neuron-selection oracle is nonconvex \\
Recursive feature machine
& Kernel metric through empirical AGOP
& Explicit iterative dimension adaptation
& Not ordinary neural training \\
Two-layer network trained by online SGD
& First-layer weights through online SGD
& Subspace convergence; single-index learning theorem
& Gaussian model, weight decay, staged training \\
\bottomrule
\end{tabular}
\caption{Different methods adapt different objects and satisfy different
guarantees; ``neural networks versus classical methods'' is too coarse.}
\label{tab:method-comparison}
\end{table}

\subsection{Entropy limits for one regression procedure}
\label{sec:yang}

The preceding results describe function classes on which particular methods
perform well. A universal entropy theorem gives a converse question that
applies to every regression procedure: if a method has risk $Cn^{-\gamma}$
at every sample size, how large can the set of targets receiving that
guarantee be?

The entropy theorem assumes a known input law $P_X$. Let
$X_1,\ldots,X_n$ be independent with common law $P_X$, and let
\begin{equation}
 Y_i=f(X_i)+\varepsilon_i,
 \qquad
 \Law(\varepsilon_1,\ldots,\varepsilon_n)
 =N(0,1)^{\otimes n},
 \label{eq:yang-model}
\end{equation}
with the noise variables independent of the inputs. A regression procedure
$\cA$ produces an estimator $\widehat f_n$ at every sample size. Define
\begin{equation}
 R_{\cA}(f,n)=\E\bigl[\|\widehat f_n-f\|_2^2\bigr]
 \label{eq:yang-risk}
\end{equation}
and, for fixed $A,C>0$ and $0<\gamma\leq1$, define its rate-$\gamma$
success set
\begin{equation}
 \cS_{\cA}(C,\gamma)
 =\left\{f:\|f\|_2\leq A,\quad
 R_{\cA}(f,n)\leq Cn^{-\gamma}
 \text{ for every }n\geq1\right\}.
 \label{eq:success-set}
\end{equation}
The requirement at every $n$ matters. A procedure can accidentally perform
well on a very large set at one fixed sample size without giving a uniform
rate along the sample-size sequence.

For a subset $\cF\subseteq L_2(P_X)$, let $M(\epsilon,\cF)$ be the logarithm
of the largest number of elements of $\cF$ separated pairwise by more than
$\epsilon$ in $L_2(P_X)$. Thus $M$ is packing entropy, not packing number.

\begin{theorem}[A universal entropy bound]
\label{thm:yang}
For every regression procedure $\cA$ and every $0<\gamma<1$, there is a
constant $C'$ depending on $A,C$, and $\gamma$ such that, for all sufficiently
small $\epsilon>0$,
\begin{equation}
 M\bigl(\epsilon,\cS_{\cA}(C,\gamma)\bigr)
 \leq C'\epsilon^{-2(1-\gamma)/\gamma}.
 \label{eq:yang-entropy}
\end{equation}
For $\gamma=1$, the corresponding upper bound has order
$\log(1/\epsilon)$.
\end{theorem}

The proof turns the sequence of estimates produced by $\cA$ into a joint
coding distribution. The cumulative risk controls the Kullback--Leibler
divergence between the coding distribution and every regression model in the
success set. A standard
testing lower bound then prevents the success set from containing too many
well-separated functions. The argument is information-theoretic and does
not depend on whether $\cA$ is a neural network, a kernel method, or a
classical smoother.

The exponent calibrates the theorem. Substituting the generic smoothness
rate
\begin{equation}
 \gamma_d=\frac{2\alpha}{2\alpha+d}
 \label{eq:smoothness-gamma}
\end{equation}
into \cref{eq:yang-entropy} gives
\begin{equation}
 \frac{2(1-\gamma_d)}{\gamma_d}=\frac d\alpha.
 \label{eq:smoothness-entropy-exponent}
\end{equation}
The packing entropy of a $d$-dimensional $\alpha$-smooth ball has the same
order $\epsilon^{-d/\alpha}$. A minimax-optimal smoothness-adaptive
procedure can therefore attain the same entropy order. Serving that entire
ball at its minimax rate does not preclude better performance on smaller
subsets.
At the slower $d$-dimensional rate, structured smooth targets may already
belong to the full smoothness ball. At a faster rate, only a smaller
structured subset receives the stronger guarantee. Even if an additional
family lies outside the ball, adding the family need not change the leading
entropy order of the success set. Saturating the entropy exponent therefore
does not exhaust a procedure's possible successes.

For an $s$-dimensional multi-index class, the analogous rate
$\gamma_s=2\alpha/(2\alpha+s)$ gives entropy exponent $s/\alpha$. Allowing
all $s$-dimensional orientations does not change the leading exponent when
$d$ is fixed. An $\epsilon$-net over the orientation matrices contributes a
finite-dimensional term of order
\begin{equation}
 O\bigl(ds\log(1/\epsilon)\bigr),
 \label{eq:orientation-entropy}
\end{equation}
which is smaller than $\epsilon^{-s/\alpha}$ as $\epsilon\downarrow0$.
The entropy theorem therefore permits one procedure to serve all orientations
and all smooth links at the $s$-dimensional rate. The theorem forbids the
same rate on a class with substantially larger packing entropy, such as the
full $d$-dimensional smoothness ball when $s<d$.

The calculation also explains why adaptation to unknown smoothness is
possible. One estimator may attain the minimax rate separately on many
nested smoothness classes. At any fixed rate, the corresponding success set
must satisfy \cref{eq:yang-entropy}; slower rates permit larger sets. The
constants need not be uniform across every class. A sharper negative result
holds for specially constructed, mutually separated classes:
no procedure can achieve each class's minimax risk with one uniform
finite-sample constant for infinitely many such classes. The theorem does
not rule out adaptation for particular countable collections.

Packing entropy measures the size of a success set but not its membership.
Two procedures can both attain the maximal entropy order while one
procedure's success set strictly contains the other's. The theorem also does
not identify which functions a given procedure favors. The entropy bound
therefore places a rate-dependent size limit on the success set of every
learning procedure, including a neural network trained by gradient descent.

\subsection{Maxisets and a concrete open question}
\label{sec:maxisets}

A \emph{maxiset} records the targets on which a fixed procedure attains a
fixed rate. For a positive sequence $r=(r_n)_{n\geq1}$, define
\begin{equation}
 \operatorname{MS}(\cA,r)
 =\left\{f:
 \sup_{n\geq1}
 \frac{\E[\|\widehat f_n-f\|_2^2]}{r_n}<\infty
 \right\}.
 \label{eq:maxiset}
\end{equation}
An ordinary risk upper bound proves that a named function class is contained
in the maxiset. An exact maxiset characterization also needs a converse:
targets outside a specified approximation or regularity space cannot attain
the rate under that procedure. Such characterizations are known for several
wavelet-thresholding and model-selection procedures.

Maxisets at different rates have a natural nesting. If
$r_n\leq C_0s_n$ for every $n$, then
\begin{equation}
 \operatorname{MS}(\cA,r)\subseteq\operatorname{MS}(\cA,s).
 \label{eq:maxiset-nesting}
\end{equation}
Faster rates have smaller success sets. A smoothness-adaptive procedure has
the relevant smoothness classes inside these nested sets. The same procedure
may also contain projection-pursuit or multi-index targets at faster rates
when those targets add sufficiently little packing entropy.

For neural networks, the important unresolved problem is algorithm-specific:
characterize the rate-indexed maxisets of a fixed architecture,
initialization, regularization rule, and gradient-training algorithm. Does
ordinary training simultaneously adapt to generic smoothness, ridge
structure, and multi-index structure? Do its success sets attain the maximal
entropy order? Do they contain the success sets of projection pursuit or a
recursive feature machine? Existing approximation theorems and restricted
feature-learning theorems do not answer these questions.

\subsection{Related feature-learning theorems}

The first theorem separates one neural feature-learning scheme from a fixed
kernel regime.
\begin{theorem}[Polynomial feature learning, informal form]
Let $\Law(X)=N(0,I_d)$ and let $f(x)=g(Ux)$ be a degree-$p$ polynomial, where
$U$ has a small number of rows. Under a nondegeneracy condition on the
expected Hessian, a two-layer training scheme learns the hidden subspace with
sample complexity of leading order $d^2$. The corresponding fixed-kernel
regime requires order $d^p$ samples.
\end{theorem}
The training scheme removes the linear component of the target, takes one
carefully scaled first-layer gradient step, resamples biases, and then trains
the output layer. The theorem proves feature learning but does not analyze
ordinary joint training.

The second theorem, discussed in the optional material of
\href{lecture13.pdf}{Lecture 13}, takes a different route.
\begin{theorem}[Intrinsic-dimensional rate from pruned gradient training, informal form]
Suppose that the input law is supported in a shrinking neighborhood of a
low-dimensional Lipschitz manifold and that the regression function is smooth
in the ambient coordinates. A prescribed randomly pruned sigmoid network,
output truncation, and full-batch gradient descent with data-dependent step
sizes achieve, up to logarithmic factors, the intrinsic-dimensional rate
\begin{equation}
 \E\bigl[\|\widehat f_n-f\|_2^2\bigr]
 \leq C(\log n)^c n^{-2p/(2p+d^*)},
 \label{eq:kohler-manifold-rate}
\end{equation}
where $p$ is the assumed smoothness and $d^*$ is the manifold dimension used
in tuning the estimator.
\end{theorem}
The theorem supplies a rate for a concrete gradient-based neural estimator.
The assumptions concern input geometry rather than hierarchical composition,
and the result does not isolate an advantage caused by depth or by feature
movement. The theorem does not compare the neural estimator with non-neural
manifold-regression methods, so the rate alone establishes no neural advantage.

\subsection{What depth changes}
\label{sec:depth}

\href{lecture14.pdf}{Lecture 14} gave an exact representation separation:
repeated composition of the tent map has a linear-size deep representation,
while a one-hidden-layer
ReLU network needs exponentially many units. Approximation theory contains two
stronger forms of the separation.
\begin{proposition}[Two depth separations, informal form]
There are oscillatory target families represented by polynomial-size deep
ReLU networks for which networks restricted to a specified smaller depth need
exponentially many units to attain a fixed approximation error. There are also
radial functions in growing dimension that three-layer networks approximate
with polynomial size, while every two-layer approximation at the stated
accuracy requires exponential width.
\end{proposition}
The proposition establishes a representation advantage for depth. The
separations do not show that gradient descent finds the deep representation
from data.

Hierarchical composition gives a statistical counterpart to the depth
separations. Consider a target of the form
\begin{equation}
 f=g_q\circ g_{q-1}\circ\cdots\circ g_0,
 \label{eq:hierarchical-composition}
\end{equation}
where each coordinate of $g_i$ depends on at most $t_i$ coordinates of its
input and has H\"older smoothness $\beta_i$. Later nonsmooth maps can reduce
the effective smoothness of an earlier constituent. Define
\begin{equation}
 \beta_i^*=\beta_i\prod_{\ell=i+1}^{q}(\beta_\ell\wedge1).
 \label{eq:effective-composition-smoothness}
\end{equation}
\begin{theorem}[Regression rate for hierarchical compositions, informal form]
Under bounded-domain and noise assumptions, an empirical risk minimizer over
a suitably sparse ReLU network class with depth of order $\log n$ satisfies
for constants $C,c>0$,
\begin{equation}
 \E\bigl[\|\widehat f_n-f\|_2^2\bigr]
 \leq C(\log n)^c
 \max_{0\leq i\leq q}
 n^{-2\beta_i^*/(2\beta_i^*+t_i)}.
 \label{eq:composition-statistical-rate}
\end{equation}
If $t_i\leq\min(d_0,\ldots,d_{i-1})$ for $i=1,\ldots,q$, the polynomial
part of the rate matches the minimax lower bound for the stated composition
class.
\end{theorem}
Each part of the network approximates one
low-dimensional constituent, and successive layers compose these
approximations. Depth therefore converts a specified compositional
description into an estimator whose rate is governed by the hardest local
constituent rather than by the ambient dimension.

The statistical conclusion has two important boundaries. The estimator is a
global empirical risk minimizer over a tuned network class; the theorem does
not show that ordinary gradient training discovers an unknown composition.
The theorem also does not prove that every shallow neural estimator must have
a slower rate. Polynomial-spline estimators can exploit related hierarchical
interaction models. The statistical result shows that deep networks can
exploit hierarchical composition. The result does not give neural networks
exclusive access to the structure.

Depth can also change the regularizer induced by parameter norm without
changing the represented function class. The relevant matrix factorization
identity makes the effect precise.
\begin{proposition}[Factorization cost]
Let $A$ have rank at most every intermediate width. For $L\geq2$,
\begin{equation}
 \inf_{A=A_L\cdots A_1}
 \frac1L\sum_{\ell=1}^L\|A_\ell\|_F^2
 =\|A\|_{S_{2/L}}^{2/L},
 \label{eq:factorization-cost}
\end{equation}
where $\|A\|_{S_p}=(\sum_j\sigma_j(A)^p)^{1/p}$ is the Schatten
$p$-quasi-norm.
\end{proposition}
For a network with several linear input layers followed by one ReLU layer,
collapsing the linear layers leaves a shallow ReLU function. Additional linear
layers therefore leave the represented function class unchanged while the
sum of squared parameter norms increasingly favors low-rank input matrices as
$L$ grows. The formal analysis connects the induced cost to single-index and
multi-index functions; experiments also study standard training. No general
statistical rate for gradient descent follows from the representation-cost
result.

The layerwise neural feature ansatz in \cref{eq:nfa} supplies a complementary
empirical claim: trained layers align with the AGOP of the remaining network.
Together, the results suggest two precise roles for depth. Nonlinear
composition can reduce representation size, and even additional linear
factorization can alter the implicit or explicit regularization. A general
theorem converting either mechanism into a statistical and computational
advantage under standard deep-network training remains unavailable.

\section{Take-home messages}

\begin{enumerate}[leftmargin=*]
\item \textbf{Approximation, constructive approximation, and learning answer
different questions.} Approximation asks whether a good function exists in
the family. Constructive approximation specifies information about the target
and an algorithm. Regression learning specifies noisy observations and a
prediction risk.

\item \textbf{Bounded integral representations have accurate finite
approximations.} If
\(f=\int\phi_\theta\,d\mu(\theta)\) in a Hilbert space and
\(\|\phi_\theta\|\leq1\), then \(m\) signed atoms attain error at most
\(\|\mu\|_{\TV}/\sqrt m\). Random sampling appears in the proof; the lemma's
conclusion is deterministic. The sampling law depends on the unknown
representing measure.

\item \textbf{Density answers the minimum approximation question.} The RBF
and ridge density results show that the error can be made arbitrarily small.
For H\"older balls, quantitative approximation theory relates error to the
number of terms; suitable neural-network constructions give corresponding
smoothness-dependent rates.

\item \textbf{Barron's theorem gives a target-dependent representation
advantage.} On a spectral Barron ball, a shallow ridge network attains an
\(m^{-1/2}\) approximation rate, while the Kolmogorov \(m\)-width for fixed
linear spaces can decay only as \(m^{-1/d}\). The comparison depends on the
target class, norm, and target-dependent choice of ridge atoms.

\item \textbf{A fixed kernel restricts the representing measure.} A fixed
feature law \(\tau\) permits measures of the form \(d\mu=p\,d\tau\) with
\(p\in L^2(\tau)\). A variation representation may use a measure concentrated
on features selected for the target. The larger representation class creates a
feature-selection problem.

\item \textbf{Gradient information can reduce a multi-index problem from
dimension \(d\) to dimension \(s\).} For
\(f(x)=g(W^\top x)\), every gradient lies in
\(\operatorname{span}(W)\). Informative gradients and an SVD recover that
subspace when gradients are available. Passive regression must first estimate
the derivative information accurately.

\item \textbf{The neural learning theorem covers a staged procedure under
restrictive assumptions.} The procedure first updates only the first-layer
weights, then freezes those weights, resamples the biases, and trains only the
output weights. For a monotone single-index model, the procedure satisfies a
prediction bound with nearly linear dependence on the ambient dimension. The
result does not show that ordinary simultaneous training realizes the general
Barron representation advantage.
\end{enumerate}
\printendnotes

\section*{Bibliographical notes}

The minimax H\"older benchmark in \cref{eq:generic-smooth-rate} is due to the
nonparametric regression theory developed by \citet{stone1982}. The infimum
there is over statistical decision rules, without a computational
restriction.

The Hilbert-space sampling argument is a standard probabilistic-method form
of finite approximation in a convex or variation ball. The wavelet route to
\cref{cor:rbf-density} follows \citet{delyon1995}: a target-dependent
continuous wavelet transform supplies the representing measure, and sampling
that measure supplies a finite scale--location approximation. The appendices
record the reconstruction and Fourier calculations needed for the Gaussian
example. Countable dyadic wavelet bases and multiresolution analysis are
developed by \citet{mallat1989}; \citet{daubechies1988} constructs compactly
supported orthonormal wavelets with high regularity. The RBF corollary uses
the continuous-wavelet construction rather than the full dyadic theory.
Classical quantitative approximation rates for H\"older and related smoothness
classes are discussed by \citet{devore1998}; \citet{petersen2026} survey
corresponding neural-network approximation results.

The shallow-network upper bound and fixed-linear-space lower bound in
\cref{thm:barron} are normalized versions of results in
\citet{barron1993}. Fourier conventions and the precise Barron norm change
constants. The anchored norm in \cref{eq:barron-norm} includes the constant
part so that the displayed quantity is a norm rather than a seminorm.
Barron's lower-bound proof uses an orthogonal family of Fourier modes. The
resulting fixed-space quantity is a Kolmogorov width. Classical nonlinear
\(m\)-term approximation and its relation to widths are surveyed by
\citet{devore1998}.

Bach~\citep{bach2017} develops the fixed-feature RKHS and target-dependent
variation-space comparison used after \cref{thm:barron}. Bach also
analyzes structured statistical rates and conditional-gradient selection from
an infinite neuron dictionary. The selection step remains a nonconvex
optimization problem for the ridge dictionary. The measure-valued formula in
\cref{eq:mean-field-representation} is standard in infinite-width two-layer
network limits; \citet{mei2019} analyze corresponding mean-field dynamics.

Projection pursuit regression was introduced by \citet{friedman1981}.
\citet{jones1992} proves the Hilbert-space greedy lemma used in
\appendixref{app:active-query-recovery} and applies the method to projection
pursuit and neural-network training.
\citet{barroncohendahmendevore2008} develop several greedy variants and their
statistical consequences. The constructive guarantees depend on the
quality and cost of the atom-selection step.

The active point-query literature makes the gradient--SVD mechanism
quantitative. \citet{cohendaubechiesdevore2012} study ridge-function recovery,
\citet{fornasierschnassvybiral2012} combine randomized queries with sparse
recovery, and \citet{tyagicevher2012} study dense multi-index and
additive-ridge models. \appendixref{app:active-query-recovery} gives an
elementary coordinate finite-difference version and an obstruction showing
why an informative-gradient condition is necessary.

Average derivative and structure-adaptive dimension-reduction methods use
related gradient information under passive noisy sampling.
\citet{hristache2001}, \citet{dalalyan2008}, and \citet{xia2002} analyze
different local-smoothing procedures and different dimension, regularity, and
identifiability regimes. The optional discussion of the AGOP, recursive
feature machines, and the neural feature ansatz follows
\citet{radhakrishnan2024}; the empirical ansatz and the available restricted
theorems must be distinguished.

The population-gradient decomposition, online contraction theorem, and
staged single-index prediction theorem are due to \citet{mousavi2023}. The
polynomial feature-learning comparison is from \citet{damian2022}. The
pruned sigmoid-network theorem is from \citet{kohler2026}. The entropy limit
and its adaptive-estimation discussion are from \citet{yang2007}; the
maxiset terminology and comparison with minimax analysis are developed by
\citet{kerkyacharian2002}.

The optional depth separations are due to \citet{telgarsky2016} and
\citet{eldan2016}. The effect of additional linear layers on representation
cost is analyzed by \citet{parkinson2025}. Hierarchical interaction models
and corresponding neural or spline estimators are studied by
\citet{kohler2017}, \citet{bauer2019}, \citet{schmidthieber2020}, and
\citet{kohlerlanger2021}. \citet{petersen2026} provide a broad reference for
modern neural-network approximation theory.
\section*{Glossary}

\begin{description}[style=unboxed,leftmargin=1.5em,labelindent=0pt]
\item[Dictionary]
A finite, countable, or continuous collection of candidate atoms used to
represent or approximate functions.
\item[Constructive approximation]
Approximation by an algorithm that receives specified information about the
target, makes finitely many oracle calls and arithmetic operations, and
returns an approximating function.
\item[Signed measure]
A real-valued measure represented in
\cref{def:signed-integral-representation} as a measurable sign function times
a finite positive measure.
\item[Total variation measure and norm]
For $d\mu=s\,d|\mu|$ with $|s|=1$, the positive measure $|\mu|$ is the
total variation measure and $\|\mu\|_{\TV}=|\mu|(\Theta)$ is the total
variation norm.
\item[H\"older continuity]
The condition $|h(x)-h(x')|\leq L\|x-x'\|^\beta$ for
$0<\beta\leq1$. Lipschitz continuity is the case $\beta=1$.
\item[Nonlinear $m$-term approximation]
Approximation by a sum of at most $m$ dictionary atoms whose identities and
coefficients may depend on the target. The selection makes the approximation
map nonlinear.
\item[Kolmogorov $m$-width]
The smallest worst-case approximation error achieved by an $m$-dimensional
linear space, as defined in \cref{eq:kolmogorov-width}.
\item[Wavelet]
A function in a translation--dilation system that satisfies a reconstruction
property. Compact support provides localization, while vanishing moments
suppress local low-degree polynomial behavior; the resulting coefficients
record both location and scale.
\item[Radial-basis-function network]
A finite sum $\sum_{j=1}^m c_j\varphi(s_j(x-t_j))$ with a fixed radial
function $\varphi$ and trainable coefficients, scales, and centers.
\item[Atomic norm]
The smallest total variation of a signed measure representing a function as
an integral of dictionary atoms. For a countable basis with a unique
coefficient expansion, the atomic norm is the $\ell_1$ norm of the
coefficients.
\item[Representation advantage]
A smaller approximation error or smaller representation size than a named
comparison class, without a claim that the representation can be estimated
from data.
\item[Learning advantage]
A risk or sample-complexity improvement for a specified estimator and data
model relative to a specified competing method.
\item[Ridge function]
A function of the form $x\mapsto g(w^\top x)$. A multi-index function uses
several linear projections.
\item[Projection pursuit]
A regression method that fits a sum of ridge functions and estimates their
directions and link functions.
\item[Atom-selection oracle]
An operation that returns a dictionary atom maximizing its correlation with
a supplied residual. Approximate variants require an additional oracle-error
term in the approximation guarantee.
\item[Active point-query model]
An observation model in which an algorithm chooses each input at which the
target function is evaluated. The model differs from passive regression,
where the input locations are sampled from a fixed distribution.
\item[Variation space]
The class of integral feature representations whose signed representing
measure has finite total variation.
\item[Spectral Barron class]
A class of restrictions of ambient Fourier representations with bounded
anchored spectral norm, as in \cref{eq:barron-norm,eq:barron-class}.
\item[Variation Barron space]
A variation space generated by a ridge-feature dictionary. Spectral Barron
conditions can imply membership in a variation Barron space, but the two
definitions are not identical. The literature uses several feature
normalizations and hence several related Barron norms.
\item[Mean-field representation]
An infinite-width two-layer network written as an integral over a probability
measure on neuron parameters. Mean-field training theory also specifies a
geometry and an evolution equation for that measure.
\item[Average gradient outer product]
The matrix $\E[\nabla f(X)\nabla f(X)^\top]$, whose leading eigenspace records
directions along which $f$ varies on average.
\item[Average derivative estimation]
A family of dimension-reduction methods that estimate averages or moments of
target derivatives from regression data in order to recover predictive
directions.
\item[Recursive feature machine]
A kernel procedure that alternates predictor fitting with an update of the
kernel metric to the fitted predictor's empirical AGOP.
\item[Neural feature ansatz]
The empirical relation in \cref{eq:nfa} between a trained layer's Gram matrix
and the AGOP of the remaining network.
\item[Feature learning]
A data-dependent change in the features or feature metric used for
prediction. The phrase alone does not specify a statistical guarantee.
\item[Hierarchical composition model]
A function $f=g_q\circ\cdots\circ g_0$ in which each coordinate of $g_i$
depends on only a bounded number of coordinates from the preceding layer.
The local arities can be much smaller than the ambient input dimension.
\item[Packing entropy]
The logarithm of the largest number of functions separated pairwise by more
than a given distance.
\item[Success set]
The functions on which a fixed procedure attains a stated risk bound at
every sample size.
\item[Maxiset]
The success set of a fixed estimation procedure at a fixed rate, usually
characterized by both an upper bound and a converse.
\end{description}

\section{Exercises}
\label{sec:exercises}

\exerciseinstructions
\setcounter{courseexercise}{0}
\exercisegroup{Integral representations and density}
\begin{exercisechunk}
\item \textbf{Monte Carlo sparsification in a Hilbert space.}
\label[exercise]{ex:l15-monte-carlo-sparsification}
\exerciseprofile{2}{2}{3}{\exproof\quad\excalculation\quad\exinterpretation}{%
Derive the $m^{-1/2}$ approximation rate and adapt the sampling law to
nonuniform bounds on the dictionary atoms.}
Let $\cH$ be a real Hilbert space, let $\|\phi_\theta\|_{\cH}\leq1$, and
suppose
\[
 f=\int\phi_\theta\,d\mu(\theta),
 \qquad B=\|\mu\|_{\TV}\in(0,\infty).
\]
Set $\pi=|\mu|/B$ and let $s=d\mu/d|\mu|$. Let
$\Theta_1,\ldots,\Theta_m$ be independent random elements with common law
$\pi$, and define
\[
 Z_j=B s(\Theta_j)\phi_{\Theta_j},
 \qquad
 \widehat f_m=\frac1m\sum_{j=1}^m Z_j.
\]
\begin{enumerate}[(a),beginpenalty=10000]
\item Prove that $\E[Z_j]=f$ and $\|Z_j\|_{\cH}\leq B$ almost surely.
\item Expand $\E[\|\widehat f_m-f\|_{\cH}^2]$ using the inner product in
$\cH$. Show that the cross terms vanish and conclude that
\[
 \E[\|\widehat f_m-f\|_{\cH}^2]\leq\frac{B^2}{m}.
\]
\item Identify the probabilistic-method step that proves the existence of a
deterministic $m$-term approximation. Then identify two pieces of information
that the construction uses but a learner observing noisy input--output pairs
does not automatically possess.
\item Suppose instead that $q:\Theta\to(0,\infty)$ is measurable,
$\|\phi_\theta\|_{\cH}\leq q(\theta)$ for $|\mu|$-almost every $\theta$, and
\[
 A=\int_\Theta q(\theta)\,d|\mu|(\theta)<\infty.
\]
Let $\widetilde\pi$ be the probability measure defined by
$d\widetilde\pi=q\,d|\mu|/A$. For independent
$\widetilde\Theta_1,\ldots,\widetilde\Theta_m$ with common law
$\widetilde\pi$, define
\[
 \widetilde Z_j
 =A\,s(\widetilde\Theta_j)
 \frac{\phi_{\widetilde\Theta_j}}{q(\widetilde\Theta_j)}.
\]
Prove that the average of the $\widetilde Z_j$ is unbiased for $f$ and has
expected squared error at most $A^2/m$. State how the original construction
is recovered as a special case.
\end{enumerate}
\end{exercisechunk}

\begin{exercisechunk}
\item \textbf{Density does not determine an approximation rate.}
\label[exercise]{ex:l15-density-versus-rate}
\exerciseprofile{2}{2}{2}{\exproof\quad\excalculation\quad\exinterpretation}{%
Distinguish density of an approximation family from the rate for an
individual target.}
Let $\cH=\ell_2(\mathbb N)$ with standard orthonormal basis
$(e_j)_{j\geq1}$, and let
$V_m=\operatorname{span}\{e_1,\ldots,e_m\}$.
\begin{enumerate}[(a),beginpenalty=10000]
\item Prove that $\bigcup_{m\geq1}V_m$ is dense in $\cH$.
\item Fix $p>0$ and define $f^{(p)}\in\cH$ by
$f_j^{(p)}=j^{-p-1/2}$. Prove that
\[
 \inf_{v\in V_m}\|f^{(p)}-v\|_2
 =\left(\sum_{j>m}j^{-2p-1}\right)^{1/2}
 \asymp m^{-p}
 \qquad\text{as }m\to\infty.
\]
\item Define $f^{(\log)}_1=0$ and
\[
 f^{(\log)}_j=\frac1{\sqrt j\log j},\qquad j\geq2.
\]
Prove that $f^{(\log)}\in\cH$ and
\[
 \inf_{v\in V_m}\|f^{(\log)}-v\|_2
 \asymp\frac1{\sqrt{\log m}}
 \qquad\text{as }m\to\infty.
\]
\item Explain which statement is common to all targets in parts (b) and (c),
and which quantitative conclusion depends on the target. Use the two examples
to state precisely why density alone does not imply a polynomial approximation
rate.
\end{enumerate}
\begin{hints}
The best approximation in $V_m$ is the orthogonal projection onto the first
$m$ coordinates. Compare each tail sum with the corresponding improper
integral.
\end{hints}
\end{exercisechunk}

\begin{exercisechunk}
\item \textbf{From H\"older approximation to regression risk.}
\label[exercise]{ex:l15-holder-bias-variance}
\exerciseprofile{2}{2}{2}{\excalculation\quad\exinterpretation}{%
Derive the smoothness-dependent regression rate by balancing approximation
and estimation errors.}
Suppose an estimator based on an $m$-dimensional approximation space satisfies
\begin{equation}
 \sup_{f\in\cH_d^\alpha(L)}
 \E\bigl[\|\widehat f_{n,m}-f\|_{L^2(P_X)}^2\bigr]
 \leq A m^{-2\alpha/d}+B\frac mn,
 \label{eq:exercise-holder-risk-bound}
\end{equation}
where $A,B>0$ do not depend on $m$ or $n$.
\begin{enumerate}[(a),beginpenalty=10000]
\item Balance the two terms in
\cref{eq:exercise-holder-risk-bound}. Show that a choice
\[
 m\asymp n^{d/(2\alpha+d)}
\]
gives risk of order $n^{-2\alpha/(2\alpha+d)}$. Explain why rounding $m$ to
a positive integer does not change the rate.
\item Suppose that the target has the form $f(x)=g(W^\top x)$, the matrix
$W\in\R^{d\times s}$ is known, and the same approximation and estimation
bounds hold in the reduced $s$-dimensional coordinates. Derive the choices
\[
 m\asymp n^{s/(2\alpha+s)},
 \qquad
 \E\bigl[\|\widehat f_{n,m}-f\|_{L^2(P_X)}^2\bigr]
 \lesssim n^{-2\alpha/(2\alpha+s)}.
\]
\item Explain why the calculation in part (b) proves neither the
$s$-dimensional rate for the full $d$-dimensional H\"older ball nor the
existence of an algorithm that finds an unknown subspace from noisy
regression data. Your explanation must identify the target-class restriction
and the missing estimation task separately.
\end{enumerate}
\end{exercisechunk}

\exercisegroup{Adaptive atoms and fixed features}
\begin{exercisechunk}
\item \textbf{Barron approximation versus a fixed linear space.}
\label[exercise]{ex:l15-barron-linear-comparison}
\exerciseprofile{3}{2}{3}{\exproof\quad\excalculation\quad\exinterpretation}{%
Prove the orthogonal-family lower-bound mechanism and compare it with
target-dependent ridge approximation.}
Let $e_1,\ldots,e_{2m}$ be orthonormal elements of a Hilbert space $\cH$,
let $V\subseteq\cH$ satisfy $\dim(V)\leq m$, and let $P_V$ be the
orthogonal projection onto $V$.
\begin{enumerate}[(a),beginpenalty=10000]
\item Prove that
\[
 \sum_{j=1}^{2m}\|P_Ve_j\|_{\cH}^2\leq m.
\]
Conclude that some $j\in\{1,\ldots,2m\}$ satisfies
\[
 \inf_{v\in V}\|e_j-v\|_{\cH}\geq\frac1{\sqrt2}.
\]
\item Let $\cF\subseteq\cH$. Suppose that, for every $m\geq1$, there are
$2m$ orthonormal elements $e_1,\ldots,e_{2m}$ such that
$\rho_m e_j\in\cF$ for every $j$, where $\rho_m>0$. Prove that the
Kolmogorov width satisfies
\[
 d_m(\cF;\cH)\geq\frac{\rho_m}{\sqrt2}.
\]
\item Assume that every $f\in\cB_C(Q)$ has an $m$-unit ridge
approximation satisfying
\[
 \|f-f_m\|_2\leq 2C m^{-1/2},
\]
and that every fixed $m$-dimensional linear space $V_m$ satisfies
\[
 \sup_{f\in\cB_C(Q)}
 \inf_{g\in V_m}\|f-g\|_2
 \geq \kappa(C/d)m^{-1/d}.
\]
Give a sufficient width for the ridge approximation to attain error
at most $\epsilon$. Give a necessary dimension for a fixed linear space to
attain worst-case error at most $\epsilon$.
\item Compare the dependence on $\epsilon$ when $d=2$ and when $d=20$.
Which choices in the ridge approximation are allowed to depend on the
target, and which choices in the linear lower bound are fixed?
\item State why the comparison proves neither a smaller minimax regression
risk for a neural-network estimator nor a contradiction with the minimax
lower bound for a full $d$-dimensional smoothness ball. Identify the missing
learning statement in the first case and the different target classes in the
second case.
\end{enumerate}
\begin{hints}
For part (a), expand $P_V$ in an orthonormal basis of $V$ and apply Bessel's
inequality to each basis vector.
\end{hints}
\end{exercisechunk}

\begin{exercisechunk}
\item \textbf{A fixed feature law can give a strictly smaller space.}
\label[exercise]{ex:l15-fixed-feature-strictness}
\exerciseprofile{2}{2}{3}{\exproof\quad\excalculation\quad\exinterpretation}{%
Use a discrete dictionary to distinguish variation representations from
fixed-feature RKHS representations.}
Let $\cH=\ell_2(\mathbb N)$, let $(e_j)_{j\geq1}$ be its standard
orthonormal basis, and set $\tau_j=2^{-j}$. Use the dictionary
$\phi_j=e_j$. Define
\begin{align*}
 \cF_1
 &=\left\{f=\sum_{j\geq1}\mu_j e_j:
                \sum_{j\geq1}|\mu_j|<\infty\right\},\\
 \cF_{2,\tau}
 &=\left\{f=\sum_{j\geq1}\tau_jp_je_j:
                \sum_{j\geq1}\tau_jp_j^2<\infty\right\},
\end{align*}
with the norms defined by the corresponding coefficient costs.
\begin{enumerate}[(a),beginpenalty=10000]
\item Use uniqueness of coordinates in the orthonormal basis to prove
\[
 \|f\|_{\cF_1}=\sum_{j\geq1}|f_j|,
 \qquad
 \|f\|_{\cF_{2,\tau}}^2
 =\sum_{j\geq1}\frac{f_j^2}{\tau_j}.
\]
\item Prove that
\[
 \|f\|_{\cF_1}\leq\|f\|_{\cF_{2,\tau}}
 \qquad\text{for every }f\in\cF_{2,\tau},
\]
and conclude that
$B_{\cF_{2,\tau}}\subseteq B_{\cF_1}$.
\item Define $f_j=2^{-j/2}/\sqrt j$. Prove that
$f\in\cF_1$ but $f\notin\cF_{2,\tau}$. Set
$g=f/\|f\|_{\cF_1}$ and prove that
$g\in B_{\cF_1}\setminus B_{\cF_{2,\tau}}$. Conclude that both inclusions
in \cref{eq:F2-F1-ball-inclusion} can be strict.
\item For the function in part (c), identify the signed-measure coefficients
$\mu_j$ and the density coefficients $p_j$. Explain why their summability
properties exhibit the restriction imposed by the fixed feature law $\tau$.
\end{enumerate}
\end{exercisechunk}

\exercisegroup{Recovering low-dimensional structure}
\begin{exercisechunk}
\item \textbf{When the AGOP recovers a multi-index subspace.}
\label[exercise]{ex:l15-agop-subspace}
\exerciseprofile{2}{2}{2}{\exproof\quad\excalculation\quad\exconstruction\quad\exinterpretation}{%
Prove the population identity behind gradient-based dimension reduction and
identify the nondegeneracy condition required for subspace recovery.}
Let $X$ be a random vector in $\R^d$, let $W\in\R^{d\times s}$ have full
column rank, and suppose $f(x)=g(W^\top x)$ with differentiable $g$. Assume
the displayed expectations below are finite.
\begin{enumerate}[(a),beginpenalty=10000]
\item Use the chain rule to prove
\[
 G_f:=\E[\nabla f(X)\nabla f(X)^\top]
 =W G_g W^\top,
 \qquad
 G_g:=\E[\nabla g(W^\top X)\nabla g(W^\top X)^\top].
\]
\item Prove that $\operatorname{range}(G_f)\subseteq
\operatorname{range}(W)$. Show that equality holds
when $G_g$ is positive definite.
\item Give an example in which $W$ has two columns but $G_f$ has rank one.
What does the example say about recovering a declared index space from
gradients? Explain why an arbitrary interpolant can have a different AGOP
from the regression target and why a recursive feature machine therefore
needs the AGOP of an accurate predictor.
\end{enumerate}
\end{exercisechunk}

\exercisegroup{Appendix-based extensions}
\begin{exercisechunk}
\item \textbf{From wavelet coefficient decay to an approximation rate.}\suppmark
\label[exercise]{ex:l15-wavelet-truncation}
\exerciseprofile{2}{2}{2}{\exproof\quad\excalculation\quad\exinterpretation}{%
Convert scale-wise wavelet coefficient bounds into an $m$-term
approximation rate.}
Let $(\psi_{j,k})$ be an orthonormal wavelet basis of $L^2([0,1])$.
Suppose that, at every level $j\geq0$, at most $C_0 2^j$ coefficients are
nonzero and
\[
 |\langle f,\psi_{j,k}\rangle_{L^2}|
 \leq C_1 2^{-j(r+1/2)},
 \qquad r>0.
\]
Let $f_J$ retain the coarse-scale terms and all wavelet coefficients from
levels $0,\ldots,J-1$.
\begin{enumerate}[(a),beginpenalty=10000]
\item Use orthonormality to prove
\[
 \|f-f_J\|_{L^2}^2
 \leq
 \frac{C_0C_1^2}{1-2^{-2r}}\,2^{-2rJ}.
\]
\item Suppose the number $m_J$ of retained coefficients satisfies
$m_J\asymp2^J$. Prove that
\[
 \|f-f_J\|_{L^2}=O(m_J^{-r}).
\]
\item State how \cref{prop:wavelet-coefficient-decay} supplies the hypothesis
of this exercise when the $r$th derivative is uniformly bounded on the
intervals appearing in that proposition.
Identify the step in parts (a)--(b) that converts local smoothness information
into a global $L^2$ approximation bound.
\end{enumerate}
\end{exercisechunk}

\begin{exercisechunk}
\item \textbf{Greedy approximation with an approximate oracle.}\suppmark
\label[exercise]{ex:l15-approximate-greedy}
\exerciseprofile{3}{2}{3}{\exproof\quad\excalculation\quad\exinterpretation}{%
Determine how accurately a correlation oracle must be solved to preserve the
$m^{-1/2}$ greedy approximation rate.}
Use the setting of \cref{prop:greedy-oracle}. Thus
$\|\phi\|_{\cH}\leq1$ for $\phi\in\cD$,
$f\in\cK=B\,\overline{\operatorname{conv}}(\cD\cup-\cD)$, and $f_0=0$.
At step $m\geq1$, write $r_{m-1}=f-f_{m-1}$ and suppose the oracle returns
$\phi_m\in\cD\cup-\cD$ satisfying
\[
 \langle r_{m-1},B\phi_m\rangle_{\cH}
 \geq
 \sup_{\phi\in\cD\cup-\cD}
 \langle r_{m-1},B\phi\rangle_{\cH}-\delta_m,
\]
where $\delta_m\geq0$. Set
\[
 f_m=\left(1-\frac1{m+1}\right)f_{m-1}
     +\frac1{m+1}B\phi_m.
\]
\begin{enumerate}[(a),beginpenalty=10000]
\item Prove by induction that $f_m\in\cK$ and hence
$\|B\phi_m-f_{m-1}\|_{\cH}\leq2B$.
\item Prove the recursion
\[
 \|r_m\|_{\cH}^2
 \leq
 \left(1-\frac{2}{m+1}\right)\|r_{m-1}\|_{\cH}^2
 +\frac{2\delta_m}{m+1}
 +\frac{4B^2}{(m+1)^2}.
\]
\item Suppose that $\delta_m\leq B^2/(m+1)$ for every $m\geq1$.
Prove by induction that
\[
 \|f-f_m\|_{\cH}^2\leq\frac{6B^2}{m+1}.
\]
\item State the required order of the oracle's additive accuracy at iteration
$m$. Explain what the result controls and what additional information would
be needed to turn the iteration bound into an arithmetic running-time bound
for a continuous ridge dictionary.
\end{enumerate}
\begin{hints}
For part (b), compare the approximate oracle value with the value of the
linear functional at $f\in\cK$.
\end{hints}
\end{exercisechunk}

\begin{exercisechunk}
\item \textbf{Stable finite differences for active subspace recovery.}\suppmark
\label[exercise]{ex:l15-noisy-finite-differences}
\exerciseprofile{2}{2}{2}{\exproof\quad\excalculation\quad\exinterpretation}{%
Quantify the tradeoff between truncation error and observation error in
finite-difference gradient estimates.}
Use the active-query setting of
\cref{thm:point-query-subspace}. Suppose
$\|\nabla^2f(x)\|_{\mathrm{op}}\leq M$ on every line segment used below,
where $M>0$. Instead of $f(x)$, the oracle returns
\[
 y(x)=f(x)+\xi(x),
 \qquad |\xi(x)|\leq\zeta,
\]
where $\zeta>0$ and the errors may be chosen adversarially. For base points
$x_1,\ldots,x_q$, define
\[
 \widetilde G_{\ell j}
 =\frac{y(x_j+he_\ell)-y(x_j)}h,
 \qquad \ell=1,\ldots,d,\quad j=1,\ldots,q.
\]
\begin{enumerate}[(a),beginpenalty=10000]
\item Prove the entrywise bound
\[
 |\widetilde G_{\ell j}-\partial_\ell f(x_j)|
 \leq\frac{Mh}{2}+\frac{2\zeta}{h}.
\]
\item Minimize the right-hand side over $h>0$. Show that
\[
 h_\star=2\sqrt{\frac\zeta M},
 \qquad
 \frac{Mh_\star}{2}+\frac{2\zeta}{h_\star}
 =2\sqrt{M\zeta}.
\]
Assume that the query domain permits the step size $h_\star$.
\item Let $G$ be the exact gradient matrix. Prove
\[
 \|\widetilde G-G\|_{\mathrm{op}}
 \leq2\sqrt{M\zeta dq}
\]
when $h=h_\star$. If $\sigma_k(G)=\lambda_\star$ and
$2\sqrt{M\zeta dq}<\lambda_\star$, use
\cref{thm:point-query-subspace} to prove
\[
 \|\widehat P-P\|_{\mathrm{op}}
 \leq
 \frac{2\sqrt{M\zeta dq}}
      {\lambda_\star-2\sqrt{M\zeta dq}}.
\]
\item Compare the noiseless limit $\zeta=0$ with the case $\zeta>0$.
Explain why decreasing $h$ improves the truncation error but can worsen the
effect of inaccurate point evaluations.
\end{enumerate}
\end{exercisechunk}

\begin{exercisechunk}
\item \textbf{Contraction toward a teacher subspace.}\suppmark
\label[exercise]{ex:l15-subspace-contraction}
\exerciseprofile{2}{2}{2}{\exproof\quad\excalculation\quad\exconstruction\quad\exinterpretation}{%
Separate contraction of irrelevant weight components from recovery and
prediction guarantees.}
Use the population-gradient decomposition in
\cref{eq:population-gradient-decomposition}. Let
$U\in\R^{s\times d}$ satisfy $UU^\top=I_s$, set
$\Pi_\perp=I_d-U^\top U$, and suppose population gradient descent satisfies
\[
 W^{(t+1)}=W^{(t)}-\eta\bigl(A_tW^{(t)}+D_tU\bigr),
\]
where $W^{(t)}\in\R^{m\times d}$, $D_t\in\R^{m\times s}$, and each
$A_t\in\R^{m\times m}$ is symmetric with
\[
 \mu I_m\preceq A_t\preceq L I_m,
 \qquad 0<\mu\leq L,\qquad 0<\eta\leq\frac1L.
\]
Define $W_\perp^{(t)}=W^{(t)}\Pi_\perp$.
\begin{enumerate}[(a),beginpenalty=10000]
\item Prove
\[
 W_\perp^{(t+1)}=(I_m-\eta A_t)W_\perp^{(t)}.
\]
Identify the algebraic property that removes the term $D_tU$.
\item Prove
\[
 \|W_\perp^{(t+1)}\|_{\mathrm F}
 \leq(1-\eta\mu)\|W_\perp^{(t)}\|_{\mathrm F}
\]
and hence
\[
 \|W_\perp^{(T)}\|_{\mathrm F}
 \leq e^{-\eta\mu T}\|W_\perp^{(0)}\|_{\mathrm F}.
\]
\item Give one matrix whose rows lie in the teacher subspace but do not span
that subspace. Use the example to explain why
$\|W_\perp^{(T)}\|_{\mathrm F}\to0$ does not prove recovery of every
teacher direction. Then state which additional network quantities and error
criterion would have to be controlled to obtain a prediction guarantee.
\end{enumerate}
\end{exercisechunk}

\exercisegroup{Further results from the literature}
\begin{exercisechunk}
\item \textbf{Balancing variation-space approximation and estimation.}\suppmark
\label[exercise]{ex:l15-bach-rate-balance}
\exerciseprofile{2}{2}{2}{\excalculation\quad\exinterpretation}{%
Recover the hidden-dimension exponent by balancing a variation-norm budget
against approximation error.}
Let
\[
 \beta=\alpha+\frac{s-1}{2},
 \qquad
 B_n(D)=A\left(\frac{D}{\sqrt n}+D^{-1/\beta}\log D\right),
 \qquad D\geq e,
\]
where $A$ is independent of $D$ and $n$.
\begin{enumerate}[(a),beginpenalty=10000]
\item Ignore the logarithmic factor and balance the two remaining terms.
Derive the resulting choice of $D$ and rate in $n$.
\item Restore the logarithm at the level of order notation and show that
the rate is
\[
 \widetilde O\left(n^{-1/(2\alpha+s+1)}\right).
\]
For $\alpha=1$, compare the exponents for a one-dimensional projection and
for a generic $d$-dimensional target.
\item Explain why replacing $d$ by $s$ is an adaptation statement rather
than a guarantee over the full $d$-dimensional smoothness class. Identify the
continuous atom-selection optimization required by a conditional-gradient
implementation and explain why the statistical analysis does not supply a
polynomial-time algorithm for that optimization.
\end{enumerate}
\end{exercisechunk}

\begin{exercisechunk}
\item \textbf{What the universal entropy bound permits.}\suppmark
\label[exercise]{ex:l15-yang-entropy}
\exerciseprofile{3}{2}{3}{\exproof\quad\excalculation\quad\exinterpretation}{%
Use the entropy exponent to distinguish permissible structural adaptation
from an impossible uniform rate on a larger class.}
Suppose a regression procedure has squared-risk success set $\cS_\gamma$ at
rate $n^{-\gamma}$. The entropy theorem gives
\[
 M(\epsilon,\cS_\gamma)
 \lesssim \epsilon^{-2(1-\gamma)/\gamma},
 \qquad 0<\gamma<1.
\]
Assume that the packing entropy of an $r$-dimensional $\alpha$-smooth ball
has order $\epsilon^{-r/\alpha}$.
\begin{enumerate}[(a),beginpenalty=10000]
\item Substitute $\gamma_r=2\alpha/(2\alpha+r)$ into the entropy exponent and
show that the upper bound matches the smoothness-ball entropy order.
\item Consider all functions $x\mapsto g(W^\top x)$, where $g$ ranges over
an $s$-dimensional $\alpha$-smooth ball and $W$ ranges over orthonormal
$d\times s$ matrices. Suppose $d$ is fixed and that covering the possible
orientations contributes $O(ds\log(1/\epsilon))$ to the entropy. Explain why
the entropy theorem permits the rate $n^{-2\alpha/(2\alpha+s)}$ uniformly over
all orientations.
\item Show why the same $s$-dimensional rate cannot hold uniformly over a
full $d$-dimensional smoothness ball when $s<d$.
\item Explain how one procedure can adapt to several unknown smoothness
levels without contradicting the entropy theorem. Your explanation should use
rate-indexed success sets and constants in addition to asymptotic notation.
\end{enumerate}
\end{exercisechunk}


\appendix
\renewcommand{\thesection}{15.\Alph{section}}

\section{A basic wavelet coefficient estimate}
\label{app:wavelet-facts}

The following one-dimensional estimate makes the role of vanishing moments
explicit. Tensor-product and multivariate wavelets satisfy analogous bounds
with multi-index derivatives.

\begin{proposition}[Decay in a smooth region]
\label{prop:wavelet-coefficient-decay}
Let $\psi\in L_1(\R)$ be compactly supported and satisfy
\begin{equation}
 \int_{\R}u^q\psi(u)\,du=0,
 \qquad q=0,\ldots,r-1.
 \label{eq:one-dimensional-moments}
\end{equation}
Define $\psi_{j,k}(x)=2^{j/2}\psi(2^jx-k)$. If $f$ is $r$ times
continuously differentiable on
$I_{j,k}=\operatorname{conv}(\{2^{-j}k\}\cup
\operatorname{supp}(\psi_{j,k}))$, then
\begin{equation}
 \bigl|\langle f,\psi_{j,k}\rangle_{L_2}\bigr|
 \leq
 \frac{2^{-j(r+1/2)}}{r!}
 \left(\int_{\R}|u|^r|\psi(u)|\,du\right)
 \sup_{x\in I_{j,k}}|f^{(r)}(x)|.
 \label{eq:wavelet-coefficient-decay}
\end{equation}
\end{proposition}

\begin{proof}
Let $x_0=2^{-j}k$, and let $P_{r-1}$ be the Taylor polynomial of $f$ at
$x_0$ through degree $r-1$. The moment conditions and the change of variable
$u=2^jx-k$ give
$\langle P_{r-1},\psi_{j,k}\rangle_{L_2}=0$. Taylor's theorem therefore
implies
\begin{align*}
 \bigl|\langle f,\psi_{j,k}\rangle_{L_2}\bigr|
 &=\bigl|\langle f-P_{r-1},\psi_{j,k}\rangle_{L_2}\bigr|\\
 &\leq
 \frac{\sup_{x\in I_{j,k}}|f^{(r)}(x)|}{r!}
 \int_{\R}|x-x_0|^r|\psi_{j,k}(x)|\,dx\\
 &=\frac{2^{-j(r+1/2)}}{r!}
 \left(\int_{\R}|u|^r|\psi(u)|\,du\right)
 \sup_{x\in I_{j,k}}|f^{(r)}(x)|.
\end{align*}
\end{proof}

The estimate is local: the derivative supremum is taken over the interval
joining the wavelet support to its translation point. Smooth regions
therefore produce rapidly decaying fine-scale coefficients, while a
singularity influences nearby wavelets. When the dyadic wavelet family forms
an orthonormal basis, Parseval's identity converts coefficient-tail bounds
into $L_2$ approximation bounds.

\section{Fourier facts used by the approximation arguments}
\label{app:fourier-basics}

The wavelet and shallow-network arguments use the same Fourier transform for
different purposes. The wavelet calculation checks a reconstruction identity
through a Fourier multiplier. The shallow-network calculation first writes a
target as an integral of Fourier modes and then writes each mode as an
integral of ridge functions. This appendix states the required Fourier facts
with the normalization fixed in \cref{eq:fourier-transform}.

\subsection{Inversion, Plancherel, and elementary identities}

A function $f:\R^d\to\mathbb C$ belongs to the \emph{Schwartz class}
$\cS(\R^d)$ if $f$ is infinitely differentiable and
\begin{equation}
 \sup_{x\in\R^d}|x^\alpha\partial^\beta f(x)|<\infty
 \quad\text{for every pair of multi-indices $\alpha,\beta$.}
 \label{eq:schwartz-class}
\end{equation}
Every smooth compactly supported function is a Schwartz function. The
changes of variables, differentiations, and integrations by parts below are
therefore justified for Schwartz functions and have no boundary terms.

\begin{proposition}[Elementary Fourier identities]
\label{prop:elementary-fourier-identities}
Let $f,g\in\cS(\R^d)$, let $t\in\R^d$, and let $s>0$. With
\begin{equation}
 (f*g)(x)=\int_{\R^d}f(y)g(x-y)\,dy,
 \label{eq:convolution-definition}
\end{equation}
the Fourier transform satisfies
\begin{align}
 \widehat{f(\,\cdot-t)}(\omega)
 &=e^{-i\omega^\top t}\widehat f(\omega),
 &
 \widehat{f(s\,\cdot)}(\omega)
 &=s^{-d}\widehat f(\omega/s),
 \label{eq:fourier-translation-dilation}\\
 \widehat{f*g}(\omega)
 &=\widehat f(\omega)\widehat g(\omega),
 &
 \widehat{\partial^\alpha f}(\omega)
 &=(i\omega)^\alpha\widehat f(\omega).
 \label{eq:fourier-convolution-differentiation}
\end{align}
In particular,
\begin{equation}
 \widehat{(-\Delta)^r f}(\omega)
 =\|\omega\|_2^{2r}\widehat f(\omega),
 \qquad r=0,1,\ldots.
 \label{eq:fourier-laplacian}
\end{equation}
If $f$ is real-valued, then
\begin{equation}
 \widehat f(-\omega)=\overline{\widehat f(\omega)}.
 \label{eq:fourier-conjugate-symmetry}
\end{equation}
\end{proposition}

\begin{proof}
The first two identities follow by the substitutions $u=x-t$ and $u=sx$.
Fubini's theorem followed by the substitution $u=x-y$ proves the convolution
identity. Integration by parts gives
$\widehat{\partial_jf}(\omega)=i\omega_j\widehat f(\omega)$; iteration proves
the multi-index formula and \cref{eq:fourier-laplacian}. Taking the complex
conjugate of \cref{eq:fourier-transform} proves
\cref{eq:fourier-conjugate-symmetry}.
\end{proof}

\begin{theorem}[Fourier inversion and Plancherel]
\label{thm:fourier-inversion-plancherel}
For every $f\in\cS(\R^d)$,
\begin{equation}
 f(x)=\frac{1}{(2\pi)^d}\int_{\R^d}
 e^{i\omega^\top x}\widehat f(\omega)\,d\omega.
 \label{eq:fourier-inversion}
\end{equation}
For every $f,g\in\cS(\R^d)$,
\begin{equation}
 \int_{\R^d}f(x)\overline{g(x)}\,dx
 =\frac{1}{(2\pi)^d}\int_{\R^d}
 \widehat f(\omega)\overline{\widehat g(\omega)}\,d\omega.
 \label{eq:fourier-parseval}
\end{equation}
Consequently, the Fourier transform extends uniquely from
$\cS(\R^d)$ to $L_2(\R^d)$ and satisfies
\begin{equation}
 \|f\|_{L_2}^2
 =\frac{1}{(2\pi)^d}\|\widehat f\|_{L_2}^2.
 \label{eq:fourier-plancherel}
\end{equation}
\end{theorem}

\begin{proof}[Proof sketch]
Let
$g_\epsilon(x)=(2\pi\epsilon)^{-d/2}
\exp(-\|x\|_2^2/(2\epsilon))$. Completing the square in the Gaussian
integral gives
$\widehat g_\epsilon(\omega)=
\exp(-\epsilon\|\omega\|_2^2/2)$. Substituting the definition of
$\widehat f$ into the right-hand side below, applying Fubini's theorem, and
evaluating the inverse Gaussian integral give
\begin{equation}
 (f*g_\epsilon)(x)
 =\frac{1}{(2\pi)^d}\int_{\R^d}e^{i\omega^\top x}
 e^{-\epsilon\|\omega\|_2^2/2}\widehat f(\omega)\,d\omega.
 \label{eq:gaussian-regularized-inversion}
\end{equation}
The functions $g_\epsilon$ have integral one and concentrate their mass in
every neighborhood of zero as $\epsilon\downarrow0$. Uniform continuity of a
Schwartz function therefore gives
$(f*g_\epsilon)(x)\to f(x)$. The right-hand side converges to
the right-hand side of \cref{eq:fourier-inversion} by dominated convergence.
Indeed, integration by parts gives
$(1+\|\omega\|_2^2)^k\widehat f(\omega)
=\widehat{(I-\Delta)^kf}(\omega)$, so $\widehat f$ is integrable for
sufficiently large $k$. Substituting the inversion formula for $g$ into
$\int f\overline g$ and applying Fubini proves
\cref{eq:fourier-parseval}. Smooth compactly supported functions are dense in
$L_2(\R^d)$ and belong to $\cS(\R^d)$, so $\cS(\R^d)$ is dense in
$L_2(\R^d)$. Continuity in \cref{eq:fourier-parseval} gives the unique extension and
\cref{eq:fourier-plancherel}.
\end{proof}

The density argument for ridge networks uses one quantitative consequence of
smoothness: a compactly supported smooth target has every polynomial spectral
moment finite.

\begin{proposition}[Rapid Fourier decay of smooth compactly supported functions]
\label{prop:smooth-fourier-decay}
If $f\in C_c^\infty(\R^d)$ and $N\geq0$, then there is a finite constant
$A_{f,N}$ such that
\begin{equation}
 |\widehat f(\omega)|
 \leq A_{f,N}(1+\|\omega\|_2)^{-N},
 \qquad \omega\in\R^d.
 \label{eq:smooth-fourier-rapid-decay}
\end{equation}
In particular,
\begin{equation}
 \int_{\R^d}\|\omega\|_1|\widehat f(\omega)|\,d\omega<\infty.
 \label{eq:first-spectral-moment-finite}
\end{equation}
\end{proposition}

\begin{proof}
For every integer $k\geq0$, \cref{eq:fourier-laplacian} gives
\begin{equation}
 (1+\|\omega\|_2^2)^k|\widehat f(\omega)|
 =|\widehat{(I-\Delta)^kf}(\omega)|
 \leq\|(I-\Delta)^kf\|_{L_1}.
 \label{eq:rapid-decay-from-derivatives}
\end{equation}
Choose $k$ with $2k\geq N$ to obtain
\cref{eq:smooth-fourier-rapid-decay}. Choosing instead
$2k>d+1$ makes
$\|\omega\|_1(1+\|\omega\|_2^2)^{-k}$ integrable and proves
\cref{eq:first-spectral-moment-finite}.
\end{proof}

\subsection{Fourier verification of the wavelet reconstruction}

Let $\psi,\varphi\in\cS(\R^d)$. For a target $f$, define the continuous
analysis transform
\begin{equation}
 u_f(s,t)=\int_{\R^d}f(x)\psi(s(x-t))\,dx,
 \qquad s>0,\quad t\in\R^d.
 \label{eq:wavelet-analysis-transform}
\end{equation}
The desired scale--location reconstruction is
\begin{equation}
 f(y)=\int_0^\infty\int_{\R^d}
 u_f(s,t)\varphi(s(y-t))s^{2d-1}\,dt\,ds
 \quad\text{in $L^2(\R^d)$}.
 \label{eq:continuous-wavelet-reconstruction}
\end{equation}
The next proposition converts \cref{eq:continuous-wavelet-reconstruction}
into a scalar condition at each nonzero frequency. Truncating the scale
integral makes the Fourier calculation ordinary; Plancherel then controls the
limit.

\begin{proposition}[Fourier criterion for a dual scale--location pair]
\label{prop:fourier-dual-wavelet-criterion}
Let $\psi,\varphi\in\cS(\R^d)$. For $f\in\cS(\R^d)$ and
$0<\epsilon<R<\infty$, define
\begin{equation}
 T_{\epsilon,R}f(y)
 =\int_\epsilon^R\int_{\R^d}
 u_f(s,t)\varphi(s(y-t))s^{2d-1}\,dt\,ds,
 \label{eq:truncated-wavelet-reconstruction}
\end{equation}
where $u_f$ is defined in \cref{eq:wavelet-analysis-transform}. Then
$T_{\epsilon,R}$ is a Fourier multiplier with
\begin{align}
 \widehat{T_{\epsilon,R}f}(\omega)
 &=M_{\epsilon,R}(\omega)\widehat f(\omega),
 \label{eq:wavelet-fourier-multiplier}\\
 M_{\epsilon,R}(\omega)
 &=\int_\epsilon^R
 \widehat\psi(-\omega/s)\widehat\varphi(\omega/s)\,\frac{ds}{s}
 =\int_{1/R}^{1/\epsilon}
 \widehat\psi(-a\omega)\widehat\varphi(a\omega)\,\frac{da}{a}.
 \label{eq:wavelet-multiplier}
\end{align}
Suppose that
\begin{equation}
 \sup_{0<\epsilon<R<\infty}
 \|M_{\epsilon,R}\|_{L_\infty(\R^d)}<\infty
 \label{eq:wavelet-uniform-multiplier-bound}
\end{equation}
and that
\begin{equation}
 \lim_{\epsilon\downarrow0,\ R\uparrow\infty}
 M_{\epsilon,R}(\omega)=1
 \quad\text{for almost every $\omega\in\R^d$.}
 \label{eq:wavelet-multiplier-limit}
\end{equation}
Then the multiplier formula extends $T_{\epsilon,R}$ to $L_2(\R^d)$ and
$T_{\epsilon,R}f\to f$ in $L_2(\R^d)$ for every $f\in L_2(\R^d)$. The limit
is the reconstruction identity in \cref{eq:continuous-wavelet-reconstruction}.
\end{proposition}

\begin{proof}
For fixed $s$, direct integration in $t$ gives
\begin{equation}
 \widehat{u_f(s,\,\cdot)}(\omega)
 =s^{-d}\widehat f(\omega)\widehat\psi(-\omega/s).
 \label{eq:wavelet-analysis-fourier-transform}
\end{equation}
The synthesis function $z\mapsto\varphi(sz)$ has Fourier transform
$s^{-d}\widehat\varphi(\omega/s)$. The convolution identity in
\cref{eq:fourier-convolution-differentiation} now gives
\cref{eq:wavelet-fourier-multiplier}. Multiplying the $s^{-d}$ factor from
the analysis transform by the $s^{-d}$ factor from the synthesis kernel
cancels the $s^{2d}$ factor in $s^{2d-1}ds$, leaving $ds/s$. The substitution
$a=1/s$ gives the second expression in \cref{eq:wavelet-multiplier}.
By Plancherel, \cref{eq:wavelet-uniform-multiplier-bound} makes
$T_{\epsilon,R}$ uniformly bounded on $\cS(\R^d)$ and hence gives a unique
bounded extension to $L_2(\R^d)$. For any $f\in L_2(\R^d)$,
\cref{eq:fourier-plancherel,eq:wavelet-multiplier-limit} and dominated
convergence then give $\|T_{\epsilon,R}f-f\|_{L_2}\to0$.
\end{proof}

\begin{proposition}[Gaussian dual wavelet pair]
\label{prop:gaussian-wavelet-verification}
Let
\begin{equation}
 \varphi(x)=e^{-\|x\|_2^2/2},
 \qquad
 \psi=c_r(-\Delta)^r\varphi,
 \qquad
 c_r=\frac{2}{(2\pi)^d\Gamma(r)},
 \label{eq:normalized-gaussian-wavelet-pair}
\end{equation}
where $r\geq1$ is an integer and $\Gamma(r)=(r-1)!$. Then the hypotheses of
\cref{prop:fourier-dual-wavelet-criterion} hold.
\end{proposition}

\begin{proof}
The Gaussian Fourier transform and \cref{eq:fourier-laplacian} give
\begin{equation}
 \widehat\varphi(\omega)=(2\pi)^{d/2}e^{-\|\omega\|_2^2/2},
 \qquad
 \widehat\psi(\omega)
 =c_r\|\omega\|_2^{2r}\widehat\varphi(\omega).
 \label{eq:gaussian-wavelet-fourier-transforms}
\end{equation}
Both transforms are real and even. For $\omega\neq0$, the full multiplier is
\begin{align}
 \int_0^\infty
 \widehat\psi(-a\omega)\widehat\varphi(a\omega)\,\frac{da}{a}
 &=c_r(2\pi)^d\int_0^\infty
 (a\|\omega\|_2)^{2r}e^{-a^2\|\omega\|_2^2}\,\frac{da}{a}\notag\\
 &=\frac{c_r(2\pi)^d}{2}\Gamma(r)=1.
 \label{eq:gaussian-wavelet-multiplier-one}
\end{align}
The truncated multipliers lie between $0$ and $1$, so they satisfy
\cref{eq:wavelet-uniform-multiplier-bound}. The $L_2$ reconstruction identity follows from
\cref{prop:fourier-dual-wavelet-criterion}.
\end{proof}

\begin{proposition}[Wavelet proof of Gaussian RBF density]
\label{prop:gaussian-rbf-density-proof}
The finite Gaussian scale--location sums in \cref{eq:rbf-density} are dense
in $L^2(\R^d)$.
\end{proposition}

\begin{proof}
Choose an integer $r>d/4$ and the pair $(\psi,\varphi)$ from
\cref{prop:gaussian-wavelet-verification}. Because
$C_c^\infty(\R^d)$ is dense in $L^2(\R^d)$, proving the claim for
$f\in C_c^\infty(\R^d)$ suffices.
For $0<s\leq1$, Fubini's theorem gives
\begin{equation}
 \int_{\R^d}|u_f(s,t)|\,dt
 \leq s^{-d}\|f\|_{L^1}\|\psi\|_{L^1}.
 \label{eq:wavelet-coarse-scale-bound}
\end{equation}
For $s\geq1$, the identity $\psi=c_r(-\Delta)^r\varphi$ and integration by
parts give
\begin{equation}
 \int_{\R^d}|u_f(s,t)|\,dt
 \leq c_rs^{-d-2r}\|(-\Delta)^rf\|_{L^1}\|\varphi\|_{L^1}.
 \label{eq:wavelet-fine-scale-bound}
\end{equation}
Because $2r>d/2$, the two estimates imply
\begin{equation}
 B_f:=\|\varphi\|_{L^2}
 \int_0^\infty\int_{\R^d}
 |u_f(s,t)|s^{3d/2-1}\,dt\,ds<\infty.
 \label{eq:wavelet-variation-norm}
\end{equation}
Normalize the synthesis atoms as
\[
 \phi_{s,t}(y)=
 \frac{s^{d/2}}{\|\varphi\|_{L^2}}\varphi(s(y-t)),
 \qquad \|\phi_{s,t}\|_{L^2}=1.
\]
The reconstruction formula in
\cref{eq:continuous-wavelet-reconstruction} is then an integral
representation by the atoms $\phi_{s,t}$ with total variation $B_f$.
Lemma~\ref{lem:mc-sparsification} supplies finite sums converging to $f$ in
$L^2$. Absorbing the normalizing factors into the coefficients gives finite
Gaussian radial-basis-function networks.
\end{proof}

\subsection{From Fourier modes to ridge functions}

Fourier inversion represents a target through the functions
$x\mapsto e^{i\omega^\top x}$. Each such function is a complex-valued ridge
function, but its univariate profile is oscillatory rather than sigmoidal.
The next identity converts the oscillatory profile into thresholds, which are
limits of steep sigmoids.

\begin{proposition}[Fourier modes as limits of sigmoidal ridge functions]
\label{prop:fourier-to-ridge}
Let $Q=[0,1]^d$, let $P_X$ be uniform on $Q$, and let $\sigma$ satisfy
\cref{eq:sigmoidal-activation}. For $\omega\in\R^d$, define
\begin{equation}
 \alpha_\omega=\sum_{j:\,\omega_j<0}\omega_j,
 \qquad
 \beta_\omega=\sum_{j:\,\omega_j>0}\omega_j.
 \label{eq:fourier-mode-range}
\end{equation}
Then $\omega^\top x\in[\alpha_\omega,\beta_\omega]$ for $x\in Q$,
$\beta_\omega-\alpha_\omega=\|\omega\|_1$, and
\begin{equation}
 e^{it}-1
 =\int_{\alpha_\omega}^{\beta_\omega}
 ie^{ib}\bigl(\I\{t\geq b\}-\I\{0\geq b\}\bigr)\,db
 \quad\text{for every $t\in[\alpha_\omega,\beta_\omega]$.}
 \label{eq:fourier-mode-threshold-identity}
\end{equation}
Moreover, for every nonzero $\omega$ and every $b\in\R$,
\begin{equation}
 \sigma\bigl(\lambda(\omega^\top x-b)\bigr)
 \longrightarrow \I\{\omega^\top x\geq b\}
 \quad\text{in $L_2(P_X)$ as $\lambda\to\infty$.}
 \label{eq:sigmoid-to-threshold}
\end{equation}
Consequently, for every nonzero $\omega$ there are finite complex-valued
networks
\begin{equation}
 g_n(x)=c_{0,n}+\sum_{j=1}^{m_n}
 c_{j,n}\sigma(a_{j,n}^\top x+b_{j,n})
 \label{eq:finite-complex-sigmoid-networks}
\end{equation}
such that
\begin{equation}
 \|g_n-(e^{i\omega^\top(\,\cdot)}-1)\|_{L_2(P_X)}\longrightarrow0,
 \qquad
 |c_{0,n}|+\sum_{j=1}^{m_n}|c_{j,n}|
 \leq2\|\omega\|_1+o(1).
 \label{eq:fourier-mode-network-limit}
\end{equation}
\end{proposition}

\begin{proof}
The formulas for $\alpha_\omega$ and $\beta_\omega$ follow by minimizing and
maximizing each coordinate contribution $\omega_jx_j$ separately. The
fundamental theorem of calculus gives
\begin{equation}
 e^{it}-1=\int_0^t ie^{ib}\,db,
 \label{eq:complex-exponential-calculus}
\end{equation}
and writing the oriented interval from $0$ to $t$ with indicators gives
\cref{eq:fourier-mode-threshold-identity}. The absolute coefficient mass of
the two indicator terms is at most
$2(\beta_\omega-\alpha_\omega)=2\|\omega\|_1$.

For nonzero $\omega$, the hyperplane
$\{x\in Q:\omega^\top x=b\}$ has Lebesgue measure zero. The two limits in
\cref{eq:sigmoidal-activation} therefore give pointwise convergence almost
everywhere in \cref{eq:sigmoid-to-threshold}; boundedness of $\sigma$ gives
$L_2(P_X)$ convergence by dominated convergence. Moreover,
\begin{equation}
 \bigl\|\I\{\omega^\top(\,\cdot)\geq b\}
 -\I\{\omega^\top(\,\cdot)\geq b'\}\bigr\|_{L_2(P_X)}^2
 =P_X\!\left(\omega^\top X\in
 [\min\{b,b'\},\max\{b,b'\})\right).
 \label{eq:threshold-l2-continuity}
\end{equation}
The probability on the right tends to zero as $b'\to b$ because
$\omega^\top X$ has no atoms.
Thus the integrand in \cref{eq:fourier-mode-threshold-identity} is a
continuous $L_2(P_X)$-valued function of $b$, and its Riemann sums converge
in $L_2(P_X)$. Replace the finitely many thresholds in each Riemann sum by
steep sigmoids using \cref{eq:sigmoid-to-threshold}. Keep the constant
indicator $\I\{0\geq b\}$ in $c_{0,n}$. The Riemann-sum and sigmoid
approximations preserve the coefficient bound up to an $o(1)$ term, proving
\cref{eq:fourier-mode-network-limit}.
\end{proof}

For a representation satisfying \cref{eq:barron-density-representation}, set
$d\nu(\omega)=(2\pi)^{-d}F(\omega)\,d\omega$. Integrating the
threshold representations in \cref{eq:fourier-mode-threshold-identity}
against $\nu$ produces a complex signed measure on the threshold ridge
parameters. Taking real parts produces a real signed ridge-parameter measure
because the represented target is real. The total variation of the real
ridge-parameter measure is at most
\begin{equation}
 2\int_{\R^d}\|\omega\|_1\,d|\nu|(\omega).
 \label{eq:fourier-to-ridge-variation-bound}
\end{equation}
The right-hand side is twice the spectral-moment term in
\cref{eq:barron-norm}. Applying \cref{lem:mc-sparsification} to the resulting
real signed measure proves the upper-bound representation used in
\cref{thm:barron}. Rapid Fourier decay from
\cref{prop:smooth-fourier-decay}, followed by approximation of an arbitrary
$L^2(Q)$ target by smooth compactly supported functions, also proves
\cref{cor:ridge-density}.

\subsection{Measure-valued spectral representations}
\label{app:fourier-measure-extension}

The density formulation in \cref{def:spectral-barron-norm} is sufficient for
the main theorem. A measure-valued formulation enlarges the class and includes
pure Fourier modes without smoothing. For a real-valued
\(f:\R^d\to\R\), let \(\mathfrak R(f)\) contain every finite complex Borel
measure \(\nu\) such that
\begin{equation}
 f(x)=f(0)+\int_{\R^d}(e^{i\omega^\top x}-1)\,d\nu(\omega)
 \quad\text{for every }x\in\R^d,
 \label{eq:barron-measure-representation}
\end{equation}
where the integral is real. Define the anchored measure cost
\begin{equation}
 \|f\|_{\cB,\mathrm{meas}}
 =|f(0)|+\inf_{\nu\in\mathfrak R(f)}
 \int_{\R^d}\|\omega\|_1\,d|\nu|(\omega).
 \label{eq:barron-measure-coefficient}
\end{equation}
The radius-(C) class on (Q) consists of restrictions of ambient functions
whose anchored measure cost is at most (C), in parallel with
\cref{eq:barron-class}. Every density
\(F\) in \cref{eq:barron-density-representation} defines
\(d\nu=(2\pi)^{-d}F\,d\omega\), so the density-defined Barron ball is
contained in the measure-defined ball with the same radius. The
Fourier-to-threshold construction in \cref{prop:fourier-to-ridge} applies
directly to \(\nu\) and gives the
same neural upper bound. The main text uses the density class because smooth
spectral approximations to cosine modes also give the fixed-space lower bound
inside that smaller class.

\subsection{Fourier modes on the cube}

The fixed-space lower bound uses the integer Fourier modes as an orthogonal
collection under the uniform law on the cube.

\begin{proposition}[Orthogonality of integer Fourier modes]
\label{prop:cube-fourier-orthogonality}
For $k,\ell\in\mathbb Z^d$,
\begin{equation}
 \int_{[0,1]^d}e^{2\pi i(k-\ell)^\top x}\,dx
 =\I\{k=\ell\}.
 \label{eq:cube-complex-fourier-orthogonality}
\end{equation}
Consequently, if $k_1,\ldots,k_N$ are nonzero and no two are equal up to
sign, then the real functions
\begin{equation}
 x\longmapsto\sqrt2\cos(2\pi k_j^\top x),
 \qquad j=1,\ldots,N,
 \label{eq:cube-cosine-orthonormal-family}
\end{equation}
are orthonormal in $L_2([0,1]^d)$.
\end{proposition}

\begin{proof}
Fubini's theorem factors the integral in
\cref{eq:cube-complex-fourier-orthogonality} into $d$ one-dimensional
integrals. For an integer $q$,
$\int_0^1e^{2\pi iqx}\,dx$ equals $1$ when $q=0$ and $0$ otherwise. The
cosine claim follows from
$2\cos u\cos v=\cos(u-v)+\cos(u+v)$.
\end{proof}

\begin{proposition}[Smooth spectral approximations to cosine modes]
\label{prop:smooth-cosine-spectra}
For every nonzero $\omega\in\R^d$ and every $0<\delta\leq1$, there is a
real-valued function $u_{\omega,\delta}$ satisfying
\cref{eq:barron-density-representation} such that
\begin{align}
 \left\|u_{\omega,\delta}
 -\sqrt2\cos(\omega^\top\,\cdot)\right\|_{L_\infty(Q)}
 &\leq\sqrt2\,\delta,
 \label{eq:smooth-cosine-approximation}\\
 \|u_{\omega,\delta}\|_{\cB(\R^d)}
 &\leq\sqrt2\bigl(1+\|\omega\|_1+\delta\bigr).
 \label{eq:smooth-cosine-spectral-moment}
\end{align}
\end{proposition}

\begin{proof}
Choose a nonnegative, even function $\rho\in C_c^\infty(\R^d)$ supported on
$\{z:\|z\|_1\leq1\}$ and satisfying
$\int_{\R^d}\rho(z)\,dz=1$. Set
$\rho_\delta(z)=\delta^{-d}\rho(z/\delta)$ and define the Fourier density
\begin{equation}
 \widehat u_{\omega,\delta}(\xi)
 =\frac{(2\pi)^d}{\sqrt2}
 \left(\rho_\delta(\xi-\omega)+\rho_\delta(\xi+\omega)\right).
 \label{eq:smoothed-cosine-density}
\end{equation}
Fourier inversion gives
\begin{equation}
 u_{\omega,\delta}(x)
 =\sqrt2\cos(\omega^\top x)
 \int_{\R^d}e^{iz^\top x}\rho_\delta(z)\,dz.
 \label{eq:smoothed-cosine-function}
\end{equation}
The last integral is real because $\rho_\delta$ is even. For $x\in Q$ and
$z$ in the support of $\rho_\delta$,
$|z^\top x|\leq\|z\|_1\|x\|_\infty\leq\delta$. Therefore
\begin{equation}
 \left|\int_{\R^d}(e^{iz^\top x}-1)\rho_\delta(z)\,dz\right|
 \leq\delta,
 \label{eq:mollified-mode-uniform-bound}
\end{equation}
which proves \cref{eq:smooth-cosine-approximation}. The triangle inequality
also gives the spectral-moment bound
\begin{align}
 \frac1{(2\pi)^d}\int_{\R^d}
 \|\xi\|_1|\widehat u_{\omega,\delta}(\xi)|\,d\xi
 &=\frac{1}{\sqrt2}\int_{\R^d}
 \left(\|\omega+z\|_1+\|-\omega+z\|_1\right)
 \rho_\delta(z)\,dz\notag\\
 &\leq\sqrt2\bigl(\|\omega\|_1+\delta\bigr).
 \label{eq:smoothed-mode-moment-calculation}
\end{align}
Because $u_{\omega,\delta}(0)=\sqrt2$, adding the anchored constant cost
proves \cref{eq:smooth-cosine-spectral-moment}.
\end{proof}

\section{Greedy and point-query recovery of ridge structure}
\label{app:active-query-recovery}

The results in this appendix separate two conditional constructions. The
first result turns the $m^{-1/2}$ convex-hull approximation bound into a
sequential greedy method, provided that a correlation oracle is available;
the result is not a near-best $m$-term oracle inequality. The second result
estimates an active linear subspace from noiseless point evaluations,
conditional on supplied base points whose gradient matrix is full rank and
well conditioned. A final example shows why smoothness and a ridge
representation alone do not guarantee such informative base points.

\subsection{Greedy approximation with a correlation oracle}

The following proposition gives the guarantee summarized in
\cref{sec:structure}. The proof also identifies the operation that can be
difficult for a continuous ridge dictionary.

\begin{proposition}[Greedy approximation with a correlation oracle]
\label{prop:greedy-oracle}
Let $\cD$ be a collection of elements of a Hilbert space $\cH$ satisfying
$\|\phi\|_{\cH}\leq1$ for every $\phi\in\cD$. Fix $B>0$, and suppose that
\begin{equation}
 f\in \cK:=B\,\overline{\operatorname{conv}}(\cD\cup-\cD).
 \label{eq:greedy-feasible-set}
\end{equation}
Set $f_0=0$. At step $m\geq1$, assume that the supremum displayed below is
attained and let a correlation oracle return
$\phi_m\in\cD\cup-\cD$ satisfying
\begin{equation}
 \langle f-f_{m-1},B\phi_m\rangle_{\cH}
 =\sup_{\phi\in\cD\cup-\cD}
   \langle f-f_{m-1},B\phi\rangle_{\cH},
 \label{eq:correlation-oracle}
\end{equation}
and update
\begin{equation}
 f_m=\left(1-\frac1{m+1}\right)f_{m-1}
     +\frac1{m+1}B\phi_m.
 \label{eq:relaxed-greedy-update}
\end{equation}
Then $f_m$ is a linear combination of at most $m$ atoms and
\begin{equation}
 \|f-f_m\|_{\cH}\leq \frac{2B}{\sqrt{m+1}}.
 \label{eq:greedy-oracle-bound}
\end{equation}
\end{proposition}

\begin{proof}
The set $\cK$ is symmetric and convex, contains $0$, and lies in the closed
ball of radius $B$. Hence $f_m\in\cK$ for every $m$. Write
$r_{m-1}=f-f_{m-1}$ and $\gamma_m=1/(m+1)$. Since $f\in\cK$, linearity and
\cref{eq:correlation-oracle} give
\begin{equation}
 \langle r_{m-1},B\phi_m\rangle_{\cH}
 \geq \langle r_{m-1},f\rangle_{\cH}.
 \label{eq:oracle-dominates-target}
\end{equation}
Consequently,
\begin{align*}
 \|f-f_m\|_{\cH}^2
 &=\|r_{m-1}-\gamma_m(B\phi_m-f_{m-1})\|_{\cH}^2\\
 &\leq (1-2\gamma_m)\|r_{m-1}\|_{\cH}^2
       +4\gamma_m^2B^2.
\end{align*}
The last term uses
$\|B\phi_m-f_{m-1}\|_{\cH}\leq2B$. Induction, starting from
$\|f-f_0\|_{\cH}\leq B$, now gives
$\|f-f_m\|_{\cH}^2\leq4B^2/(m+1)$.
The update adds at most one new atom at every step, so $f_m$ uses at most
$m$ atoms.
\end{proof}

The proposition concerns the number of selected atoms. The computational
content of the proposition depends on \cref{eq:correlation-oracle}. The oracle is elementary for
an explicitly listed finite dictionary. For the ridge dictionary in
\cref{eq:projection-pursuit-model}, the oracle must choose both a direction
and a univariate profile. Discretizing all directions can require a set whose
size is exponential in $d$ at a fixed angular resolution. The rate in
\cref{eq:greedy-oracle-bound} therefore does not by itself control the
arithmetic cost of an iteration.

\subsection{Recovering a multi-index subspace from point values}

Noiseless active queries permit a different construction when the target is
known to have a multi-index form. Write
\begin{equation}
 \cB_d(R)=\{x\in\R^d:\|x\|_2\leq R\}.
 \label{eq:euclidean-ball-notation}
\end{equation}
The next theorem uses coordinate finite differences so that the conditional
subspace-estimation algorithm and its cost are explicit. The theorem assumes
that informative base points have already been supplied.

\begin{theorem}[Active-subspace recovery from noiseless point queries]
\label{thm:point-query-subspace}
Let $f:\cB_d(1+\bar h)\to\R$ satisfy
\begin{equation}
 f(x)=g(Ax),
 \qquad A\in\R^{k\times d},\quad AA^\top=I_k,
 \label{eq:appendix-multi-index-model}
\end{equation}
where $g$ is twice continuously differentiable. Define
\begin{equation}
 P=A^\top A,
 \quad
 L=\sup_{z\in\cB_d(1+\bar h)}\|\nabla f(z)\|_2,
 \quad
 M=\sup_{z\in\cB_d(1+\bar h)}\|\nabla^2f(z)\|_{\mathrm{op}}.
 \label{eq:active-query-regularity}
\end{equation}
Choose $x_1,\ldots,x_q\in\cB_d(1)$ and form
\begin{equation}
 G=\begin{bmatrix}\nabla f(x_1)&\cdots&\nabla f(x_q)\end{bmatrix}.
 \label{eq:sampled-gradient-matrix}
\end{equation}
Assume that $\operatorname{rank}(G)=k$ and write
$\lambda_\star=\sigma_k(G)>0$. For $0<h\leq\bar h$, use the
$q(d+1)$ point values at $x_j$ and $x_j+he_\ell$ to define
\begin{equation}
 \widetilde G_{\ell j}
 =\frac{f(x_j+he_\ell)-f(x_j)}{h},
 \qquad \ell=1,\ldots,d,\quad j=1,\ldots,q.
 \label{eq:coordinate-finite-difference-matrix}
\end{equation}
Let $\widehat P$ be the orthogonal projector onto the leading $k$ left
singular vectors of $\widetilde G$, and set
\begin{equation}
 \eta=\frac{Mh\sqrt{dq}}2.
 \label{eq:finite-difference-perturbation}
\end{equation}
If $\eta<\lambda_\star$, then
\begin{equation}
 \|\widehat P-P\|_{\mathrm{op}}
 \leq \frac{\eta}{\lambda_\star-\eta}.
 \label{eq:projector-recovery-bound}
\end{equation}
The reduced function $\widehat f(x)=f(\widehat Px)$ consequently satisfies
\begin{equation}
 \sup_{x\in\cB_d(1)}|f(x)-\widehat f(x)|
 \leq L\frac{\eta}{\lambda_\star-\eta}.
 \label{eq:active-query-function-bound}
\end{equation}
\end{theorem}

\begin{proof}
The chain rule gives
\begin{equation}
 \nabla f(x)=A^\top\nabla g(Ax).
 \label{eq:appendix-gradient-factorization}
\end{equation}
Each column of $G$ is therefore a linear combination of the rows of $A$.
The columns of $G$ span a $k$-dimensional space by the rank assumption, and
the $k$ rows of $A$ are orthonormal. Consequently, the columns of $G$ and the
rows of $A$ span the same subspace of $\R^d$. The matrix $P=A^\top A$ is the
orthogonal projector onto that subspace.

Taylor's theorem along the coordinate direction $e_\ell$ gives
\begin{equation}
 \left|\widetilde G_{\ell j}-\partial_\ell f(x_j)\right|
 \leq \frac{Mh}{2}.
 \label{eq:coordinate-difference-error}
\end{equation}
For $E=\widetilde G-G$, summing the squared entrywise bounds yields
\begin{equation}
 \|E\|_{\mathrm{op}}
 \leq\|E\|_{\mathrm F}
 \leq\frac{Mh\sqrt{dq}}2=\eta.
 \label{eq:gradient-matrix-perturbation}
\end{equation}

Let $\widehat U\widehat\Sigma\widehat V^\top$ contain the leading $k$
singular triplets of $\widetilde G$. The singular-value perturbation
inequality gives
$\sigma_k(\widetilde G)\geq\lambda_\star-\eta>0$. Since
$(I-P)G=0$,
\begin{equation}
 (I-P)\widehat U
 =(I-P)E\widehat V\widehat\Sigma^{-1}.
 \label{eq:subspace-perturbation-identity}
\end{equation}
Therefore
\begin{equation}
 \|(I-P)\widehat U\|_{\mathrm{op}}
 \leq\frac{\eta}{\lambda_\star-\eta}.
 \label{eq:one-sided-subspace-bound}
\end{equation}
The two subspaces have the same dimension, so the left-hand side of
\cref{eq:one-sided-subspace-bound} equals
$\|\widehat P-P\|_{\mathrm{op}}$. The one-sided subspace bound therefore
proves \cref{eq:projector-recovery-bound}.

Finally, $AP=A$ implies $f(x)=f(Px)$. Both $Px$ and $\widehat Px$ belong to
$\cB_d(1)$ when $x\in\cB_d(1)$. The gradient bound in
\cref{eq:active-query-regularity} and the mean-value theorem give
\[
 |f(x)-f(\widehat Px)|
 =|f(Px)-f(\widehat Px)|
 \leq L\|(P-\widehat P)x\|_2,
\]
which proves \cref{eq:active-query-function-bound}.
\end{proof}

If $\widehat U$ has orthonormal columns spanning the estimated subspace, then
\begin{equation}
 \widehat f(x)=\widehat g(\widehat U^\top x),
 \qquad
 \widehat g(y)=f(\widehat Uy),\qquad y\in\cB_k(1).
 \label{eq:recovered-link}
\end{equation}
Each query of $\widehat g$ requires one additional query of $f$. Suppose that
a grid, spline, or another $k$-dimensional approximation method uses such
queries to construct $\widetilde g$ satisfying
$\sup_{y\in\cB_k(1)}|\widehat g(y)-\widetilde g(y)|\leq\delta$. Then
\begin{equation}
 \widetilde f(x)=\widetilde g(\widehat U^\top x)
 \quad\Longrightarrow\quad
 \sup_{x\in\cB_d(1)}|f(x)-\widetilde f(x)|
 \leq L\frac{\eta}{\lambda_\star-\eta}+\delta.
 \label{eq:subspace-plus-link-error}
\end{equation}
The additional query count is the cost of obtaining the $k$-dimensional
interpolation error $\delta$. For a generic smooth link, the interpolation cost can depend
exponentially on the intrinsic dimension $k$, but not on the ambient dimension
$d$.

The rank assumption in \cref{thm:point-query-subspace} is the finite-sample
form of an average-gradient condition. If a random element $X$ has a declared
law supported on $\cB_d(1)$, define
\begin{equation}
 H_f=\E[\nabla f(X)\nabla f(X)^\top]
 =A^\top\E[\nabla g(AX)\nabla g(AX)^\top]A.
 \label{eq:active-query-population-agop}
\end{equation}
A positive $k$th eigenvalue means that the gradients reveal every active
direction on average. Strict positivity gives identifiability; a useful
finite-sample guarantee requires a quantitative lower bound on that eigenvalue.
Randomized recovery results combine such a lower bound with matrix
concentration to choose enough base points. For sparse or compressible
gradients, the directional differences
\begin{equation}
 \frac{f(x_j+h\varphi_i)-f(x_j)}h
 \approx \varphi_i^\top\nabla f(x_j)
 \label{eq:compressed-directional-differences}
\end{equation}
form a compressed-sensing problem. Sparse recovery followed by the same SVD
can then use fewer than $d$ directions per base point.

\subsection{Why an informative-link condition is necessary}

The next example shows that active queries do not make every smooth ridge
class tractable. The obstruction is informational: most queries return zero,
so the algorithm receives no evidence about the hidden direction.

\begin{proposition}[A point-query obstruction for smooth ridge functions]
\label{prop:ridge-query-lower-bound}
There are numerical constants $c_0,c_1>0$ such that the following statement
holds for every sufficiently large $d$. Define
\begin{equation}
 g(t)=8(t-1/2)_+^3,
 \qquad
 f_a(x)=g(a^\top x),\quad a\in\mathbb S^{d-1},\quad x\in\cB_d(1).
 \label{eq:hidden-cap-ridge-functions}
\end{equation}
The link $g$ is twice continuously differentiable and
$\|f_a\|_\infty=1$. Any deterministic adaptive algorithm using at most
$c_0e^{c_1d}$ noiseless point queries in $\cB_d(1)$ has the same query
transcript on the zero function and on some $f_a$. Consequently, the algorithm
cannot guarantee uniform error strictly smaller than $1/2$ on
\begin{equation}
 \{0\}\cup\{f_a:a\in\mathbb S^{d-1}\}.
 \label{eq:hidden-cap-function-class}
\end{equation}
\end{proposition}

\begin{proof}
Let $U$ be uniform on $\mathbb S^{d-1}$. Rotational invariance and the density
of the first coordinate of $U$ imply that, for every fixed
$x\in\cB_d(1)$,
\begin{equation}
 \P(U^\top x>1/2)
 \leq C\sqrt d\,(3/4)^{(d-3)/2}
 \leq C_0e^{-c_1d}
 \label{eq:spherical-cap-bound}
\end{equation}
for numerical constants $C,C_0,c_1>0$. If $x=0$, the probability is zero.
Otherwise, rotate $x/\|x\|_2$ to the first coordinate. Since
$\|x\|_2\leq1$, the event in \cref{eq:spherical-cap-bound} implies that the
first coordinate of $U$ exceeds $1/2$. The density of the first coordinate is proportional to
$(1-t^2)^{(d-3)/2}$ on $[-1,1]$. Integrating the first-coordinate density
gives the first inequality.

Run the algorithm along the transcript in which every query returns zero.
The resulting query points $x_1,\ldots,x_q$ are fixed. If
$qC_0e^{-c_1d}<1$, the union bound applied to
\cref{eq:spherical-cap-bound} gives a direction $a$ satisfying
$a^\top x_j\leq1/2$ for every $j$. Hence $f_a(x_j)=0$ for every query, and
the algorithm produces the same output, say $h$, for the targets $0$ and
$f_a$. Since $f_a(a)=1$,
\[
 1=\|f_a-0\|_\infty
 \leq\|f_a-h\|_\infty+\|h-0\|_\infty.
\]
At least one of the two errors is at least $1/2$. Reducing $c_0$ if necessary
ensures $qC_0e^{-c_1d}<1$ for every
$q\leq c_0e^{c_1d}$.
\end{proof}

The lower bound does not conflict with
\cref{thm:point-query-subspace}. Fix $a\in\mathbb S^{d-1}$ and let the random
element $X$ be uniform on $\mathbb S^{d-1}$. For the target $f_a$ in
\cref{eq:hidden-cap-ridge-functions},
\begin{equation}
 0<\lambda_1(H_{f_a})
 =\E\!\left[g'(a^\top X)^2\right]
 \leq36\P(a^\top X>1/2)
 \leq36C_0e^{-c_1d}.
 \label{eq:hidden-cap-small-agop}
\end{equation}
Thus the informative eigenvalue is positive but exponentially small. The
recovery guarantee in \cref{thm:point-query-subspace} is useful only when the
target supplies gradient information
at base points that the query strategy can obtain with the desired cost.

\section{Feature representations of reproducing-kernel Hilbert spaces}
\label{app:rkhs-features}

The next construction characterizes kernels that admit the feature
representation in \cref{eq:feature-kernel} and identifies changes of feature
law that preserve the RKHS norm.

\begin{proposition}[Every separable RKHS has a scalar feature representation]
\label{prop:every-rkhs-feature-representation}
Let $\cH_k$ be a separable real reproducing-kernel Hilbert space with kernel
$k$, and let $(e_j)_{j\in J}$ be an orthonormal basis of $\cH_k$, where $J$
is finite or countably infinite. Choose numbers $\tau_j>0$ satisfying
$\sum_{j\in J}\tau_j=1$, and define
\begin{equation}
 \phi_j(x)=\frac{e_j(x)}{\sqrt{\tau_j}}.
 \label{eq:rkhs-feature-construction}
\end{equation}
Then
\begin{equation}
 k(x,x')=\sum_{j\in J}\tau_j\phi_j(x)\phi_j(x'),
 \label{eq:rkhs-kernel-expansion}
\end{equation}
and every $f\in\cH_k$ has a representation
\begin{equation}
 f(x)=\sum_{j\in J}\tau_jp_j\phi_j(x),
 \qquad
 \|f\|_{\cH_k}^2=\sum_{j\in J}\tau_jp_j^2.
 \label{eq:rkhs-norm-expansion}
\end{equation}
\end{proposition}

\begin{proof}
For $x\in\cX$, the reproducing property gives
$e_j(x)=\langle e_j,k(x,\cdot)\rangle_{\cH_k}$. Expanding both kernel
sections in the orthonormal basis gives
\[
 k(x,x')
 =\langle k(x,\cdot),k(x',\cdot)\rangle_{\cH_k}
 =\sum_{j\in J}e_j(x)e_j(x'),
\]
which becomes \cref{eq:rkhs-kernel-expansion} after substituting
\cref{eq:rkhs-feature-construction}. If
$f=\sum_jc_je_j$, set $p_j=c_j/\sqrt{\tau_j}$. Then
\[
 \sum_j\tau_jp_j\phi_j(x)=\sum_jc_je_j(x)=f(x),
 \qquad
 \sum_j\tau_jp_j^2=\sum_jc_j^2=\|f\|_{\cH_k}^2.
\]
\end{proof}

Different bases and weights give different feature representations of the
same kernel. More generally, suppose $d\widetilde\tau=q\,d\tau$, where
$q>0$ $\tau$-almost everywhere and $\int q\,d\tau=1$. The simultaneous
changes
\begin{equation}
 \widetilde\phi_v=\frac{\phi_v}{\sqrt q(v)},
 \qquad
 \widetilde p(v)=\frac{p(v)}{\sqrt q(v)}
 \label{eq:feature-measure-change}
\end{equation}
preserve the kernel, the represented function, and the coefficient norm:
\begin{align*}
 \int\widetilde\phi_v(x)\widetilde\phi_v(x')\,d\widetilde\tau(v)
 &=\int\phi_v(x)\phi_v(x')\,d\tau(v),\\
 \int\widetilde p(v)\widetilde\phi_v(x)\,d\widetilde\tau(v)
 &=\int p(v)\phi_v(x)\,d\tau(v),\\
 \int|\widetilde p(v)|^2\,d\widetilde\tau(v)
 &=\int|p(v)|^2\,d\tau(v).
\end{align*}
Changing $\tau$ while holding the features fixed generally changes the
kernel. For a redundant fixed feature representation, several coefficient
functions $p$ may represent the same $f$; the RKHS norm is the infimum of
$\|p\|_{L_2(\tau)}$ over those coefficient functions.

\section{First-stage contraction for the staged-SGD result}
\label{app:sgd-contraction}

The prediction theorem in \cref{sec:mousavi} uses a specialized single-index
model. The first-stage analysis applies more generally to a multiple-index
response. Let \(\Law(X)=N(0,I_d)\), let the noise \(\varepsilon\) be independent
of \(X\), and suppose
\begin{equation}
 Y=g(UX;\varepsilon),
 \qquad U\in\R^{s\times d},\quad UU^\top=I_s,
 \label{eq:mousavi-model}
\end{equation}
where \(z\mapsto g(z;\varepsilon)\) is weakly differentiable for every fixed
noise value. The rows of \(U\) span the teacher subspace. The student is the
two-layer network in \cref{eq:two-layer-student}.

The contraction theorem has two assumption regimes. The smooth-activation
regime uses squared loss, an almost-sure bound on \(Y\), and uniform bounds on
\(\sigma\), \(\sigma'\), and \(\sigma''\). The ReLU regime instead assumes
\begin{equation}
 0\leq\partial_{11}^2\ell\leq1,
 \qquad |\partial_1\ell|\leq1,
 \qquad |\partial_{12}^2\ell|\leq1.
 \label{eq:mousavi-loss-assumptions}
\end{equation}
Appropriately scaled Huber and logistic losses satisfy the ReLU conditions.
Monotonicity of the link is not required for the contraction theorem.

The key identity comes from Gaussian integration by parts.
\begin{proposition}[Population-gradient decomposition, informal form]
Under the smooth-activation assumptions, or under the corresponding
conditions for a ReLU network and a convex Lipschitz loss, the gradient of the
regularized population risk has the form
\begin{equation}
 \nabla_W R_\lambda(W)
   =\bigl(H(W)+\lambda I_m\bigr)W+D(W)U,
 \label{eq:population-gradient-decomposition}
\end{equation}
where \(H(W)\in\R^{m\times m}\) and \(D(W)\in\R^{m\times s}\) depend on the
current network and the data distribution. Let \(\Pi_\perp\) be the
orthogonal projection onto the complement of the row space of \(U\), and
write \(W_\perp=W\Pi_\perp\). Right-multiplying
\cref{eq:population-gradient-decomposition} by \(\Pi_\perp\) removes the target
term:
\begin{equation}
 \nabla_W R_\lambda(W)\Pi_\perp
 =\bigl(H(W)+\lambda I_m\bigr)W_\perp.
 \label{eq:orthogonal-gradient}
\end{equation}
\end{proposition}
If the symmetric part of \(H(W)+\lambda I_m\) remains uniformly positive
along the trajectory and the step size is sufficiently small,
\cref{eq:orthogonal-gradient} contracts \(W_\perp\). Only the component inside
the teacher subspace receives the forcing term \(D(W)U\).

\begin{theorem}[Contraction toward the teacher subspace by online SGD, informal form]
The estimator updates \(W\) according to
\cref{eq:mousavi-first-stage}, using a fresh independent pair
\((X_t,Y_t)\) at every update. At initialization, the entries of \(W^{(0)}\)
are independent \(N(0,1/d)\) variables, the coordinates of \(a^{(0)}\) are
independent and uniform on \([-1/m,1/m]\), and the coordinates of \(b^{(0)}\)
are independent uniform signs. The output weights \(a^{(0)}\) and biases
\(b^{(0)}\) remain fixed. For the specified loss, weight decay, and step-size
schedules, the decaying-step version satisfies, for \(0<\delta<1\), with
probability at least \(1-\delta\),
\begin{equation}
 \frac{\|W^{(T)}_\perp\|_F}{\sqrt m}
 \ \lesssim\
 \sqrt{\frac{d+\log(1/\delta)}{T}},
 \label{eq:sgd-subspace-rate}
\end{equation}
up to constants determined by the regularization schedule. The guarantee
requires \(m\gtrsim\log(1/\delta)\).
\end{theorem}

The theorem controls only the component orthogonal to the teacher subspace.
The conclusion allows \(W^{(T)}\) to approach the zero matrix and does not
show that the rows of \(W^{(T)}\) span every teacher direction. The theorem
also gives neither convergence of the component inside the subspace nor
accurate prediction for every link \(g\).

\begin{thebibliography}{99}
\raggedright
\small

\bibitem[Bach(2017)]{bach2017}
Francis Bach.
\newblock Breaking the curse of dimensionality with convex neural networks.
\newblock \emph{Journal of Machine Learning Research}, 18(19):1--53, 2017.
\newblock \url{https://jmlr.org/papers/v18/14-546.html}.

\bibitem[Barron(1993)]{barron1993}
Andrew R. Barron.
\newblock Universal approximation bounds for superpositions of a sigmoidal
function.
\newblock \emph{IEEE Transactions on Information Theory}, 39(3):930--945,
1993.
\newblock \url{https://doi.org/10.1109/18.256500}.

\bibitem[Barron et al.(2008)]{barroncohendahmendevore2008}
Andrew R. Barron, Albert Cohen, Wolfgang Dahmen, and Ronald A. DeVore.
\newblock Approximation and learning by greedy algorithms.
\newblock \emph{The Annals of Statistics}, 36(1):64--94, 2008.
\newblock \url{https://doi.org/10.1214/009053607000000631}.

\bibitem[Bauer and Kohler(2019)]{bauer2019}
Benedikt Bauer and Michael Kohler.
\newblock On deep learning as a remedy for the curse of dimensionality in
nonparametric regression.
\newblock \emph{The Annals of Statistics}, 47(4):2261--2285, 2019.
\newblock \url{https://doi.org/10.1214/18-AOS1747}.

\bibitem[Cohen et al.(2012)]{cohendaubechiesdevore2012}
Albert Cohen, Ingrid Daubechies, Ronald DeVore, G\'erard Kerkyacharian, and
Dominique Picard.
\newblock Capturing ridge functions in high dimensions from point queries.
\newblock \emph{Constructive Approximation}, 35(2):225--243, 2012.
\newblock \url{https://doi.org/10.1007/s00365-011-9147-6}.

\bibitem[Dalalyan et al.(2008)]{dalalyan2008}
Arnak S. Dalalyan, Anatoly Juditsky, and Vladimir Spokoiny.
\newblock A new algorithm for estimating the effective dimension-reduction
subspace.
\newblock \emph{Journal of Machine Learning Research}, 9:1647--1678, 2008.
\newblock \url{https://jmlr.org/papers/v9/dalalyan08a.html}.

\bibitem[Damian et al.(2022)]{damian2022}
Alex Damian, Jason D. Lee, and Mahdi Soltanolkotabi.
\newblock Neural networks can learn representations with gradient descent.
\newblock In \emph{Proceedings of the 35th Conference on Learning Theory},
volume 178 of \emph{Proceedings of Machine Learning Research}, pages
5413--5452, 2022.
\newblock \url{https://proceedings.mlr.press/v178/damian22a.html}.

\bibitem[Daubechies(1988)]{daubechies1988}
Ingrid Daubechies.
\newblock Orthonormal bases of compactly supported wavelets.
\newblock \emph{Communications on Pure and Applied Mathematics},
41(7):909--996, 1988.
\newblock \url{https://doi.org/10.1002/cpa.3160410705}.

\bibitem[Delyon et al.(1995)]{delyon1995}
Bernard Delyon, Anatoli Juditsky, and Albert Benveniste.
\newblock Accuracy analysis for wavelet approximations.
\newblock \emph{IEEE Transactions on Neural Networks}, 6(2):332--348, 1995.
\newblock \url{https://doi.org/10.1109/72.363469}.

\bibitem[DeVore(1998)]{devore1998}
Ronald A. DeVore.
\newblock Nonlinear approximation.
\newblock \emph{Acta Numerica}, 7:51--150, 1998.
\newblock \url{https://doi.org/10.1017/S0962492900002816}.

\bibitem[Eldan and Shamir(2016)]{eldan2016}
Ronen Eldan and Ohad Shamir.
\newblock The power of depth for feedforward neural networks.
\newblock In \emph{Proceedings of the 29th Conference on Learning Theory},
volume 49 of \emph{Proceedings of Machine Learning Research}, pages
907--940, 2016.
\newblock \url{https://proceedings.mlr.press/v49/eldan16.html}.

\bibitem[Fornasier et al.(2012)]{fornasierschnassvybiral2012}
Massimo Fornasier, Karin Schnass, and Jan Vyb\'iral.
\newblock Learning functions of few arbitrary linear parameters in high
dimensions.
\newblock \emph{Foundations of Computational Mathematics}, 12(2):229--262,
2012.
\newblock \url{https://doi.org/10.1007/s10208-012-9115-y}.

\bibitem[Friedman and Stuetzle(1981)]{friedman1981}
Jerome H. Friedman and Werner Stuetzle.
\newblock Projection pursuit regression.
\newblock \emph{Journal of the American Statistical Association},
76(376):817--823, 1981.
\newblock \url{https://doi.org/10.1080/01621459.1981.10477729}.

\bibitem[Hristache et al.(2001)]{hristache2001}
Marian Hristache, Anatoli Juditsky, J\"org Polzehl, and Vladimir Spokoiny.
\newblock Structure adaptive approach for dimension reduction.
\newblock \emph{The Annals of Statistics}, 29(6):1537--1566, 2001.
\newblock \url{https://doi.org/10.1214/aos/1015345954}.

\bibitem[Jones(1992)]{jones1992}
Lee K. Jones.
\newblock A simple lemma on greedy approximation in Hilbert space and
convergence rates for projection pursuit regression and neural network
training.
\newblock \emph{The Annals of Statistics}, 20(1):608--613, 1992.
\newblock \url{https://doi.org/10.1214/aos/1176348546}.

\bibitem[Kerkyacharian and Picard(2002)]{kerkyacharian2002}
G\'erard Kerkyacharian and Dominique Picard.
\newblock Minimax or maxisets?
\newblock \emph{Bernoulli}, 8(2):219--253, 2002.
\newblock \url{https://projecteuclid.org/journals/bernoulli/volume-8/issue-2/Minimax-or-maxisets/10.3150/bj/1068128976.full}.

\bibitem[Kohler and Krzy\.{z}ak(2017)]{kohler2017}
Michael Kohler and Adam Krzy\.{z}ak.
\newblock Nonparametric regression based on hierarchical interaction models.
\newblock \emph{IEEE Transactions on Information Theory}, 63(3):1620--1630,
2017.
\newblock \url{https://doi.org/10.1109/TIT.2016.2634401}.

\bibitem[Kohler and Langer(2021)]{kohlerlanger2021}
Michael Kohler and Sophie Langer.
\newblock On the rate of convergence of fully connected deep neural network
regression estimates.
\newblock \emph{The Annals of Statistics}, 49(4):2231--2249, 2021.
\newblock \url{https://doi.org/10.1214/20-AOS2034}.

\bibitem[Kohler et al.(2026)]{kohler2026}
Michael Kohler, Adam Krzy\.{z}ak, and Vincent Molinero R\"omer.
\newblock Learning of deep neural network regression estimates using gradient
descent with pruning.
\newblock 2026.
\newblock \url{https://arxiv.org/abs/2608.25743}.

\bibitem[Mallat(1989)]{mallat1989}
St\'ephane G. Mallat.
\newblock A theory for multiresolution signal decomposition: The wavelet
representation.
\newblock \emph{IEEE Transactions on Pattern Analysis and Machine
Intelligence}, 11(7):674--693, 1989.
\newblock \url{https://doi.org/10.1109/34.192463}.

\bibitem[Mei et al.(2019)]{mei2019}
Song Mei, Theodor Misiakiewicz, and Andrea Montanari.
\newblock Mean-field theory of two-layers neural networks: Dimension-free
bounds and kernel limit.
\newblock In \emph{Proceedings of the 32nd Conference on Learning Theory},
volume 99 of \emph{Proceedings of Machine Learning Research}, pages
2388--2464, 2019.
\newblock \url{https://proceedings.mlr.press/v99/mei19a.html}.

\bibitem[Mousavi-Hosseini et al.(2023)]{mousavi2023}
Alireza Mousavi-Hosseini, Sejun Park, Manuela Girotti, Ioannis Mitliagkas,
and Murat A. Erdogdu.
\newblock Neural networks efficiently learn low-dimensional representations
with SGD.
\newblock 2023.
\newblock \url{https://arxiv.org/abs/2209.14863}.

\bibitem[Parkinson et al.(2025)]{parkinson2025}
Suzanna Parkinson, Greg Ongie, and Rebecca Willett.
\newblock ReLU neural networks with linear layers are biased towards single-
and multi-index models.
\newblock \emph{SIAM Journal on Mathematics of Data Science}, 7(3):1021--1052,
2025.
\newblock \url{https://doi.org/10.1137/24M1672158}.

\bibitem[Petersen and Zech(2026)]{petersen2026}
Philipp Petersen and Jakob Zech.
\newblock Mathematical theory of deep learning.
\newblock 2026.
\newblock \url{https://arxiv.org/abs/2407.18384}.

\bibitem[Radhakrishnan et al.(2024)]{radhakrishnan2024}
Adityanarayanan Radhakrishnan, Daniel Beaglehole, Parthe Pandit, and
Mikhail Belkin.
\newblock Mechanism for feature learning in neural networks and
backpropagation-free machine learning models.
\newblock \emph{Science}, 383(6690):1461--1467, 2024.
\newblock \url{https://doi.org/10.1126/science.adi5639}.

\bibitem[Schmidt-Hieber(2020)]{schmidthieber2020}
Johannes Schmidt-Hieber.
\newblock Nonparametric regression using deep neural networks with ReLU
activation function.
\newblock \emph{The Annals of Statistics}, 48(4):1875--1897, 2020.
\newblock \url{https://doi.org/10.1214/19-AOS1875}.

\bibitem[Stone(1982)]{stone1982}
Charles J. Stone.
\newblock Optimal global rates of convergence for nonparametric regression.
\newblock \emph{The Annals of Statistics}, 10(4):1040--1053, 1982.
\newblock \url{https://doi.org/10.1214/aos/1176345969}.

\bibitem[Telgarsky(2016)]{telgarsky2016}
Matus Telgarsky.
\newblock Benefits of depth in neural networks.
\newblock In \emph{Proceedings of the 29th Conference on Learning Theory},
volume 49 of \emph{Proceedings of Machine Learning Research}, pages
1517--1539, 2016.
\newblock \url{https://proceedings.mlr.press/v49/telgarsky16.html}.

\bibitem[Tyagi and Cevher(2012)]{tyagicevher2012}
Hemant Tyagi and Volkan Cevher.
\newblock Active learning of multi-index function models.
\newblock In \emph{Advances in Neural Information Processing Systems 25},
pages 1475--1483, 2012.
\newblock \url{https://papers.nips.cc/paper/4820-active-learning-of-multi-index-function-models}.

\bibitem[Xia et al.(2002)]{xia2002}
Yingcun Xia, Howell Tong, W. K. Li, and Li-Xing Zhu.
\newblock An adaptive estimation of dimension reduction space.
\newblock \emph{Journal of the Royal Statistical Society: Series B},
64(3):363--410, 2002.
\newblock \url{https://www.stat.purdue.edu/~yuzhu/stat598m3/Papers/AdaptiveEstimation.pdf}.

\bibitem[Yang(2007)]{yang2007}
Yuhong Yang.
\newblock How powerful can any regression learning procedure be?
\newblock In \emph{Proceedings of the 11th International Conference on
Artificial Intelligence and Statistics}, volume 2 of \emph{Proceedings of
Machine Learning Research}, pages 636--643, 2007.
\newblock \url{https://proceedings.mlr.press/v2/yang07a.html}.

\end{thebibliography}

\end{document}
